\documentclass[11pt]{article}
\usepackage[margin=1in]{geometry}
\usepackage{amsmath,amssymb,amsthm}
\usepackage{hyperref}
\usepackage{listings}
\usepackage{xcolor}
\usepackage{booktabs}
\usepackage{array}
\usepackage{microtype}
\usepackage{cite}
\usepackage{url}
\usepackage{fancyhdr}
\usepackage{multirow}
\usepackage{pdflscape}
\usepackage{rotating}

% Allow line breaks at underscores and dots in \texttt
\renewcommand{\texttt}[1]{%
  \begingroup\ttfamily\hyphenchar\font=`\-\relax
  \def\_{\textunderscore\discretionary{}{}{}}%
  #1\endgroup}

% Global stretch to prevent overfull boxes from long identifiers
\emergencystretch=3em

\newtheorem{theorem}{Theorem}
\newtheorem{lemma}[theorem]{Lemma}
\newtheorem{proposition}[theorem]{Proposition}
\newtheorem{corollary}[theorem]{Corollary}
\theoremstyle{definition}
\newtheorem{definition}[theorem]{Definition}
\newtheorem{finding}[theorem]{Finding}
\theoremstyle{remark}
\newtheorem{remark}[theorem]{Remark}

\lstset{
  basicstyle=\ttfamily\small,
  keywordstyle=\color{blue},
  commentstyle=\color{gray},
  breaklines=true,
  frame=single,
  xleftmargin=1em,
  escapeinside={(*}{*)},
}

\title{Project Yanasse: A Hofstadterian Strange Loop\\at Mathematics' Friction Points\\[0.5em]
\large Part~11: Summary of the Initial Campaign\\[0.3em]
\normalsize 100 Transfer Attempts, 10 Papers, 35 New Theorems}
\author{Alexandre Linhares\thanks{alexandre@linhares.ltd. Affiliations: Getulio Vargas Foundation and ARGO LABS}}
\date{April 2026}

\pagestyle{fancy}
\fancyhf{}
\renewcommand{\headrulewidth}{0.4pt}
\fancyhead[L]{\footnotesize\textit{Citation: Linhares, A. ``Project Yanasse: A Hofstadterian Strange Loop at Mathematics' Friction Points --- Part~11,'' Argolab.org Report for Dissemination, 2026.}}
\fancyfoot[C]{\thepage}

\begin{document}
\maketitle

\begin{abstract}
This paper summarizes the initial Yanasse campaign: 100 schema
transfer attempts across 10 (source, target) area pairs, reported
individually in Parts 1--10.  The system discovers new proofs by
transferring tactic invocation patterns from one area of Lean 4's
Mathlib library to a structurally distant area, using a
domain-independent relational analogy engine (Deep Vision) to
identify structurally similar proof states across areas, and an AI
reasoning agent (Claude) to semantically adapt the source tactic's
mathematical strategy to the target theorem.
%
Of 100 attempts, 35 produced Lean-verified new proofs (35\%
overall).  The results reveal a two-factor model of
transferability: success depends on (1) the source area's tactic
profile (structural combinators transfer broadly; algebraic decision
procedures transfer narrowly) and (2) the target area's type-class
landscape (areas with rich algebraic instances are accessible to
more sources).  An element/morphism paradigm boundary separates
areas where algebraic tactics gain purchase from those where they
cannot.  Information Theory is the most accessible target (90\%);
Condensed Mathematics the least (10\%).  Probability is the
stronger source (46\%); Number Theory the weaker (10\%), but the
two access non-overlapping parts of shared targets.
%
We distill lessons for improving the general theorem prover: which
areas to prioritize as analogy sources, which QAP score ranges are
predictive, which tactic classes transfer reliably, and where the
matching engine's blind spots lie.
\end{abstract}

\newpage
\tableofcontents
\newpage

%======================================================================
\section{Introduction}\label{sec:intro}
%======================================================================

Can a proof technique routine in one branch of mathematics be
systematically discovered in a distant branch where it is unknown?
The Yanasse project is a large-scale empirical test of this
question, using relational analogy as the discovery mechanism.

The project takes its name from Horacio Hideki Yanasse, who taught
that the best ideas cross disciplinary boundaries.  The underlying
engine---Deep Vision~\cite{linhares_deep_vision,
linhares2024deepvision}---grows out of the Fluid Analogies Research
Group tradition~\cite{hofstadter_fluid_1995, hofstadter_go_1999,
hofstadter_surfaces_2013}, which holds that analogy-making---perceiving
shared relational structure across superficially different
situations---is the core of cognition.  Deep Vision represents any
structured domain as a \emph{relational network}: entities connected
by typed, directed relations.  An analogy optimization finds the
entity permutation that maximizes relational consistency between two
networks.  The same matching code that discovers chess
analogies~\cite{linhares_active_2005, linhares_understanding_2007,
linhares_questioning_2010, linhares_entanglement_2012,
linhares_what_2013, linhares_emergence_2014} discovers proof-state
analogies---only the relation extractor is domain-specific.

The proof-state matcher has previously been used to discover a formal
proof of Wolstenholme's theorem~\cite{linhares_wolstenholme_2026} in
Lean~4 (see also~\cite{linhares_sketch_2026}).

The Yanasse campaign applies this framework at scale.  We extract
tactic usage distributions across 27~top-level Mathlib namespaces
(217,133~proof states), compute $z$-scores to identify tactics heavily
used in a \emph{source} area but rare or absent in a \emph{target}
area, match source and target proof states via GPU-accelerated analogy
on Apple's MPS backend, and then ask an AI reasoning agent to
\emph{semantically adapt}---not symbol-substitute---the source
tactic's invocation pattern to the target theorem.  The adapted proof
is compiled in Lean~4 with full Mathlib; exit code~0 with no
\texttt{sorry} constitutes verification.

This summary paper synthesizes results from Parts~1
through~10~\cite{yanasse_part1, yanasse_part2, yanasse_part3,
yanasse_part4, yanasse_part5, yanasse_part6, yanasse_part7,
yanasse_part8, yanasse_part9, yanasse_part10}, covering 2~source
areas (Probability and Number Theory), 8~target areas, 100~schema
transfer attempts, and 35~new Lean-verified proofs.

%======================================================================
\section{Related Work}\label{sec:related}
%======================================================================

The Yanasse project sits at the intersection of several active
research areas: neural and AI-assisted theorem proving, computational
analogy, and the large-scale formalization of mathematics.  We
survey each in turn, positioning the present contribution.

\paragraph{Neural theorem proving.}
The application of machine learning to formal theorem proving has
accelerated dramatically since 2020.  Polu and
Sutskever~\cite{polu2020gptf} showed that a GPT-based language model
could generate novel proofs accepted into the Metamath library,
establishing the ``generative tactic prediction'' paradigm.  Han
et~al.~\cite{han2022pact} extended this with Proof Artifact
Co-Training (PACT), using self-supervised data from Lean proof terms
to improve tactic prediction from 32\% to 48\%.
Bansal et~al.~\cite{bansal2019holist} introduced HOList, a
reinforcement-learning environment for higher-order logic that
benchmarked deep proof search at scale.
Yang et~al.~\cite{yang2024leandojo} built LeanDojo, coupling
retrieval-augmented premise selection (ReProver) with Lean~4's
Mathlib to outperform GPT-4 zero-shot.
Song et~al.~\cite{song2024leancopilot} integrated LLM inference
natively into Lean as Lean~Copilot, automating 74\% of proof steps
in interactive use.
Xin et~al.~\cite{xin2024deepseekprover} fine-tuned DeepSeek-Math on
8~million synthetic Lean~4 proofs, reaching 52\% on miniF2F.
Wang et~al.~\cite{wang2024legoprover} introduced LEGO-Prover, which
grows a reusable lemma library during proving and advanced miniF2F
from 48\% to 57\%.
Trinh et~al.~\cite{trinh2024alphageometry} demonstrated
AlphaGeometry, a neuro-symbolic prover that solved 25 of 30 Olympiad
geometry problems---matching a human gold medalist---by synthesizing
millions of training theorems.
The AlphaProof system~\cite{alphaproof2025nature} combined
reinforcement learning with Lean~4 to achieve silver-medal
performance at IMO~2024 and was published in \emph{Nature} in 2025.
Li et~al.~\cite{li2024survey} provide a comprehensive survey of
the field.

These systems all learn \emph{within} a single formal library.
Yanasse differs in two respects: it transfers tactics
\emph{across} distant areas of the same library, and its matching
engine is not a neural network but a domain-independent
relational analogy optimizer (QAP).  The analogy engine sees
proof-state \emph{shape}, not token sequences, making its
successes and failures interpretable in terms of mathematical
structure---particularly the type-class barriers that neural
provers also encounter but rarely diagnose.

\paragraph{Premise selection and hammers.}
A complementary line of work focuses on selecting relevant lemmas
from large libraries.  Kaliszyk and
Urban~\cite{kaliszyk2014flyspeck} showed that machine-learned
premise selection could prove 39\% of the Flyspeck library
automatically, establishing the ``hammer'' paradigm.  Blanchette
et~al.~\cite{blanchette2016hammering} surveyed hammer methods
across multiple proof assistants.  Gauthier, Kaliszyk, and
Urban~\cite{gauthier2021tactictoe} implemented TacticToe, which
learns from human proofs to guide Monte Carlo tree search at the
tactic level, proving 66\% of HOL4's standard library.
Mikula et~al.~\cite{mikula2024magnushammer} replaced
Sledgehammer's symbolic retrieval with a transformer-based
contrastive model (Magnushammer), outperforming it on PISA and
miniF2F.  Yanasse's analogy engine performs a function analogous
to premise selection---identifying structurally relevant source
proofs---but matches entire proof states rather than individual
lemmas, and the retrieved information is a tactic
\emph{invocation pattern}, not a premise set.

\paragraph{Cross-library alignment and proof transfer.}
The problem of transferring mathematical knowledge across formal
systems has received attention in the proof-repair and
concept-alignment communities.  Gauthier and
Kaliszyk~\cite{gauthier2019aligning} aligned concepts across six
proof assistant libraries using disambiguation on structural
similarity.  Ringer et~al.~\cite{ringer2021proofrepair}
developed Pumpkin Pi for proof repair across type equivalences in
Coq, showing that proofs can be systematically transported when the
underlying types change.  These efforts transfer knowledge across
\emph{libraries} or \emph{type representations}; Yanasse transfers
tactic strategies across \emph{mathematical domains} within a single
library, which to our knowledge has not been attempted at this scale.

\paragraph{Computational analogy.}
Deep Vision grows out of the Fluid Analogies Research Group
(FARG) tradition~\cite{hofstadter_fluid_1995, hofstadter_go_1999,
hofstadter_surfaces_2013}, which models analogy as the perception
of shared relational structure---not a post-hoc comparison, but the
core act of cognition.  Gentner's structure-mapping
theory~\cite{gentner1983structure} formalized one family of
constraints on analogical mapping; the FARG tradition emphasizes a
different set of priorities: active perception, context-sensitive
re-representation, and the inseparability of analogy from concept
formation.  Mitchell~\cite{mitchell2021analogy} surveyed the gap
between current AI and human-level abstraction and analogy-making,
noting that no existing system combines the flexibility of symbolic
structure mapping with the scalability of neural methods.
Deep Vision bridges this gap for the proof domain:
it uses a \emph{quadratic assignment} formulation and runs on GPU
hardware, enabling the matching of 217{,}133 proof states that the
Yanasse campaign requires.

The ``strange loop'' structure of the Yanasse project---where the
analogy engine's failure to prove the paradigm-boundary theorem is
itself an instance of that boundary---resonates with Hofstadter's
concept of tangled hierarchies of
self-reference~\cite{hofstadter_go_1999, hofstadter2007loop}.

\paragraph{Formalization at scale.}
The feasibility of the Yanasse campaign depends on the maturity of
large formal libraries.  The Mathlib community~\cite{mathlib2020}
has built a library of over 189{,}000 theorems in
Lean~4~\cite{moura2021lean4}, covering a broad range of
undergraduate and graduate mathematics.  Buzzard, Commelin, and
Massot~\cite{buzzard2020perfectoid} formalized perfectoid spaces,
demonstrating that Lean can handle frontier algebraic geometry.
The Liquid Tensor Experiment~\cite{scholze2022lte} verified
a theorem in condensed mathematics, one of the target areas in
our campaign.  Avigad~\cite{avigad2024formalturn} reflects on how
formalization and AI are jointly reshaping mathematical practice.
Zheng et~al.~\cite{zheng2022minif2f} introduced the miniF2F
cross-system benchmark, and Wu
et~al.~\cite{wu2022autoformalization} showed that LLMs can
autoformalize competition mathematics.  Jiang
et~al.~\cite{jiang2023dsp} combined informal reasoning with formal
proof sketches in the Draft, Sketch, and Prove framework.
Szegedy~\cite{szegedy2020autoformalization} argued that
autoformalization is a promising path toward general mathematical
reasoning.  The scale and diversity of Mathlib is what makes
cross-area tactic transfer both possible and interesting: without
$>$200{,}000 formalized proof states spanning 27~areas, the
structural map of transferability (Section~\ref{sec:map}) could not
be drawn.

%======================================================================
\section{Goals of the Experiment}\label{sec:goals}
%======================================================================

The Yanasse campaign was designed to answer four questions:

\begin{enumerate}
\item \textbf{Can relational analogy discover cross-domain proof
  techniques?}
  The core hypothesis is that a domain-independent analogy
  engine---one that sees only the \emph{shape} of proof states, not
  their mathematical content---can identify structural connections
  between distant areas of formalized mathematics that are
  sufficient to guide tactic transfer.

\item \textbf{What determines whether a tactic transfers?}
  Tactic transfer requires both a source pattern and a receptive
  target.  We sought to identify the structural factors on each side
  that predict success or failure, beyond the trivial observation
  that ``similar areas are easier.''

\item \textbf{Are there systematic barriers to transfer?}
  If certain kinds of tactics \emph{never} transfer to certain kinds
  of targets, understanding why would be as valuable as the
  successful transfers.  We sought to map these barriers empirically.

\item \textbf{Can different source areas complement each other?}
  If two source areas access different layers of the same target,
  a multi-source campaign could achieve broader coverage than any
  single source.  We tested this by sending both Probability and
  Number Theory to overlapping targets.
\end{enumerate}

The campaign was calibrated against a pinned target
(\texttt{nlinarith [sq\_nonneg ...]} from Analysis to Combinatorics)
before scaling.  The calibration itself failed---the pinned schema
was already present in the target area---but the pipeline was
validated on an alternative pair, and the failure informed the
schema-filtering rules used throughout the campaign.

%======================================================================
\section{Summary of Results}\label{sec:results}
%======================================================================

\subsection{Cumulative results across all ten pairs}

Table~\ref{tab:cumulative} presents the complete results.  Across
100~schema transfer attempts on 10~area pairs, 35~produced
Lean-verified new proofs.

\begin{table}[ht]
\centering\small
\caption{Cumulative results across all ten Yanasse pairs.}
\label{tab:cumulative}
\begin{tabular}{@{}rlcccl@{}}
\toprule
\# & Pair & Schemas & New thm & Rate & Pool$^*$ \\
\midrule
1 & Prob $\to$ RepTheory
  & 10 & 4 & 40\% & 672 \\
2 & Prob $\to$ SetTheory
  & 10 & 4 & 40\% & 2{,}076 \\
3 & Prob $\to$ Condensed
  & 10 & 2 & 20\% & 189 \\
4 & Prob $\to$ FieldTheory
  & 10 & 4 & 40\% & 2{,}792 \\
5 & Prob $\to$ Dynamics
  & 10 & 4 & 40\% & 495 \\
6 & Prob $\to$ ModelTheory
  & 10 & 5 & 50\% & 1{,}079 \\
7 & Prob $\to$ InfoTheory
  & 10 & 9 & 90\% & 240 \\
8 & NumTh $\to$ Condensed
  & 10 & 0 & 0\% & 189 \\
9 & NumTh $\to$ RepTheory
  & 10 & 0 & 0\% & 672 \\
10 & NumTh $\to$ Dynamics
  & 10 & 3 & 30\% & 495 \\
\midrule
& \textbf{Total} & \textbf{100} & \textbf{35} & \textbf{35\%}
  & --- \\
\bottomrule
\end{tabular}

\smallskip\noindent{\footnotesize $^*$Pool = number of proof states
in the target area's Mathlib corpus, against which each source proof
is matched.}
\end{table}

\subsection{Results by source area}

\begin{table}[ht]
\centering\small
\caption{Transfer rates by source area.}
\label{tab:by_source}
\begin{tabular}{@{}lcccc@{}}
\toprule
Source & Pairs & Schemas & Successes & Rate \\
\midrule
Probability & 7 & 70 & 32 & 46\% \\
Number Theory & 3 & 30 & 3 & 10\% \\
\midrule
\textbf{Total} & \textbf{10} & \textbf{100} & \textbf{35}
  & \textbf{35\%} \\
\bottomrule
\end{tabular}
\end{table}

Probability's 46\% rate reflects the domain-agnostic character of
its distinctive tactics: structural combinators
(\texttt{filter\_upwards}, \texttt{ext1}, \texttt{congr with},
\texttt{any\_goals}, \texttt{by\_cases}) that operate on logical
structure rather than algebraic content.  Number Theory's 10\% rate
reflects the opposite: its distinctive tactics
(\texttt{ring}, \texttt{decide}, \texttt{norm\_num},
\texttt{linarith}, \texttt{field}) are algebraic decision procedures
that require specific type-class instances to execute.

\subsection{Results by target area}

\begin{table}[ht]
\centering\small
\caption{Transfer rates by target area (all sources combined).}
\label{tab:by_target}
\begin{tabular}{@{}lccc@{}}
\toprule
Target area & Attempts & Successes & Rate \\
\midrule
Information Theory & 10 & 9 & 90\% \\
Model Theory & 10 & 5 & 50\% \\
Field Theory & 10 & 4 & 40\% \\
Set Theory & 10 & 4 & 40\% \\
Dynamics & 20 & 7 & 35\% \\
Representation Theory & 20 & 4 & 20\% \\
Condensed & 20 & 2 & 10\% \\
\bottomrule
\end{tabular}
\end{table}

The target difficulty spectrum ranges from 90\% (Information
Theory) to 10\% (Condensed Mathematics).  Notably, targets tested
with both source areas (Dynamics, Representation Theory, Condensed)
show that adding Number Theory as a second source adds new theorems
only for Dynamics---the one target among the three that contains
element-level arithmetic sub-layers.

\subsection{Discussion of results}

The 35\% overall rate means that roughly one in three attempts at
transferring a proof technique across Mathlib area boundaries
produces a Lean-verified new proof.  This is substantially above
what a random-guess baseline would yield: the schemas were selected
precisely because they are \emph{rare or absent} in the target area,
so the prior probability that they would compile at all is low.

The variance across pairs is large (0\%--90\%), and this variance
is itself the most informative result of the campaign.  It reveals
that transfer success is not a property of the method alone---it
depends on a structured interaction between source tactic profiles
and target type-class landscapes, as we analyze in the following
sections.

%======================================================================
\section{What Transfers}\label{sec:transfers}
%======================================================================

Across the 35~successful transfers, several patterns emerge.

\subsection{Structural combinators transfer broadly}

The most reliable transfers come from Probability's
\emph{structural combinators}---tactics whose function is to
reorganize proof state rather than to compute with algebraic
content:

\begin{itemize}
\item \texttt{ext1} (function/morphism extensionality):
  succeeded in 5 of 7~target areas (RepTheory, SetTheory,
  Condensed, Dynamics, InfoTheory).  Extensionality is
  mathematically universal: every area has objects whose equality
  reduces to pointwise equality.

\item \texttt{by\_cases} (case splitting on a structural property):
  succeeded in 5 of 7~target areas (RepTheory, SetTheory,
  FieldTheory, ModelTheory, InfoTheory).  Every area has natural
  dichotomies (zero/nonzero, empty/nonempty, finite/infinite).

\item \texttt{any\_goals} (dispatch any trivially closeable
  sub-goal): succeeded in 5 of 7~target areas (RepTheory,
  SetTheory, FieldTheory, ModelTheory, InfoTheory).  After a
  case-split or decomposition, at least one branch is often trivial.

\item \texttt{congr with} / \texttt{congrm} (congruence with
  named variables): succeeded in 3--4 targets.  Peeling one
  structural layer and naming the bound variable is a
  domain-independent operation.

\item \texttt{push} (\texttt{push\_neg}, \texttt{push\_cast},
  membership normalization): succeeded in 3 of 7~targets.
  Negation pushing and membership normalization operate on
  logical connectives, not domain-specific content.
\end{itemize}

\begin{finding}[Structural combinators as universal transfer vectors]%
\label{find:structural}
Probability's structural combinators---\texttt{ext1},
\texttt{by\_cases}, \texttt{any\_goals}, \texttt{congr with},
\texttt{push}---transfer to 5 or more of 7~target areas.  Their
universality comes from operating on logical structure
(quantifiers, disjunctions, equality, membership) rather than
algebraic content.  These are the most reliable class of tactics
for cross-area transfer.
\end{finding}

\subsection{The (head, modifier) decomposition}

Part~1 identified that tactic schemas decompose into a
\emph{head} (the tactic name) and a \emph{modifier} (the argument
structure: \texttt{with}-clause, bracketed lemma list, arity
pattern).  Across the campaign, this decomposition proved to be the
single most predictive structural feature:

\begin{itemize}
\item The \textbf{head} is typically \emph{domain-gated}: its
  elaborator requires domain-specific goal shapes or type-class
  instances.  \texttt{filter\_upwards} requires
  \texttt{Filter.Eventually}; \texttt{ring} requires
  \texttt{CommSemiring}; \texttt{measurability} requires
  $\sigma$-algebra structure.

\item The \textbf{modifier} often expresses a \emph{domain-general
  proof pattern}: ``introduce a fresh element and name it,''
  ``coordinate over a list of facts,'' ``dispatch any trivial
  sub-goal.''  These patterns have analogs across areas.
\end{itemize}

When the head fails but the modifier carries independent content,
semantic adaptation can replace the head with a target-specific
tactic while preserving the modifier's proof pattern.  This is the
mechanism behind the most interesting transfers: \texttt{filter\_upwards
[LIST] with $\omega$} becomes \texttt{ext1 $\omega$ + simp [LIST]}
in Representation Theory (Part~1); the same pattern recurs in Set
Theory, Model Theory, and Information Theory.

\subsection{Algebraic tactics transfer narrowly but genuinely}

Number Theory's algebraic decision procedures---\texttt{ring},
\texttt{ring\_nf}, \texttt{linarith}, \texttt{norm\_num}---succeeded
only 3~times in 30~attempts, all against Dynamics.  But when they
succeed, the transfer is genuinely mathematical: \texttt{ring}
closes an algebraic identity in the translation-number computation
for negative integer powers (Part~10, item~90); \texttt{ring\_nf}
normalizes an arithmetic sub-step in a semiconjugacy proof
(Part~10, item~92).  These are not structural repackaging---they
are computational steps that the target proofs actually need.

\subsection{Meta-tactics transfer independently of domain content}

Part~10 revealed that \texttt{enter}---a tactic for navigating into
nested expressions---succeeds with a QAP score of 0.0.  Its utility
depends on expression \emph{nesting depth}, not on mathematical
content.  The relational analogy matcher, which compares
entity-relation networks, cannot detect this kind of
technique-level similarity.  \texttt{enter} transferred to Dynamics
(Part~10) despite zero structural similarity between the Number
Theory and Dynamics proof states, because both domains contain
expressions nested under scalar multiplication or function
application where a surgical inner rewrite is needed.

\subsection{Automation substitution as a transfer mode}

Several successful transfers are \emph{horizontal substitutions}:
one automation tactic replaces another that accomplishes the same
mathematical step.  \texttt{measurability} replaces
\texttt{fun\_prop} in Information Theory (Part~7, item~65);
\texttt{all\_goals} replaces manual per-goal dispatch in Field
Theory (Part~4) and Model Theory (Part~6).  These transfers are
less novel than genuine technique transfers, but they are real: the
adapted proofs compile and close goals that the original
automation did not reach in the target area.

\subsection{Convergence points signal deep structural affinity}

When multiple schemas independently route to the same target
theorem via the QAP matcher, the cause matters.  In Information
Theory (Part~7), five schemas converge on
\texttt{integral\_llr\_compProd\_eq\_add}---the chain rule for
log-likelihood ratios---producing five distinct new proofs of the
same theorem.  Each proof articulates the same mathematical step
differently (\texttt{congr with}, \texttt{ext1},
\texttt{measurability}, \texttt{congrm}, \texttt{split\_ifs}).
This \emph{convergence point} signals that the source and target
share deep mathematical infrastructure, and the QAP matcher
correctly identifies it.

In contrast, the \emph{degenerate attractor} phenomenon (Part~6:
six schemas converge on \texttt{instSubsingleton}, a trivial
3-entity target) signals pool scarcity, not structural affinity.
Distinguishing genuine convergence from degenerate attraction is
diagnostically important: entity count and schema diversity are the
key discriminators.

%======================================================================
\section{What Does Not Transfer}\label{sec:barriers}
%======================================================================

The 65~failures are at least as informative as the 35~successes.
They cluster into three systematic barrier types.

\subsection{The type-class barrier}

The most common failure mode: the tactic requires a type-class
instance (\texttt{CommSemiring}, \texttt{Field},
\texttt{LinearOrderedCommSemiring}, \texttt{Decidable},
\texttt{CanLift}) that the target's proof goals do not carry.
This barrier operates at Lean~4's type-class synthesis
level---before any mathematical reasoning begins, the tactic's
elaborator fails to find the required instance.

This barrier is absolute when the target area operates entirely at
the morphism level.  In Condensed Mathematics (Part~8), proof goals
involve functors, natural transformations, and adjunctions---types
carrying no algebraic structure.  In Representation Theory (Part~9),
proof goals involve module homomorphisms and chain-complex
morphisms---types carrying additive group structure but not ring
structure.  All 20~Number Theory attempts against these two targets
failed for this reason.

\begin{finding}[The element/morphism paradigm boundary]%
\label{find:paradigm}
There exists a paradigm boundary in Mathlib between areas whose
proofs operate on \emph{elements} of algebraic structures (rings,
fields, ordered semirings) and areas whose proofs operate on
\emph{morphisms} between categorical objects (functors, natural
transformations, module homomorphisms).  Algebraic decision
procedures can cross this boundary only where the target contains
element-level sub-layers---as in Dynamics' translation numbers and
Birkhoff averages (Part~10), which embed $\mathbb{R}$-arithmetic
within otherwise topological proofs.
\end{finding}

\subsection{The goal-shape barrier}

Some tactics require a specific goal shape that the target area
simply does not generate.  \texttt{fun\_prop} requires goals tagged
with \texttt{@[fun\_prop]} attributes (\texttt{Continuous},
\texttt{Measurable}, \texttt{Differentiable}).
\texttt{filter\_upwards} requires \texttt{Filter.Eventually} or
filter membership goals.  \texttt{measurability} requires
$\sigma$-algebra predicates.  When the target area has zero such
goals, no amount of semantic adaptation can bridge the gap.

This barrier differs from the type-class barrier in that it is
about the \emph{structure of the goal proposition} rather than the
\emph{algebraic structure of the types involved}.
\texttt{fun\_prop} failed in 5 of 7~Probability targets and in all
3~Number Theory targets---making it the least transferable tactic
in the campaign.

\subsection{The homogeneity barrier}

\texttt{all\_goals} requires every sub-goal after a decomposition
to be closeable by a single decision procedure.  In areas where
decomposition produces \emph{heterogeneous} sub-goals (a
homology-class equality alongside a cycle-membership proof, or a
measurability side-condition alongside an arithmetic identity),
\texttt{all\_goals} fails even though the individual sub-goals are
tractable.  Its sister tactic \texttt{any\_goals}, which tolerates
heterogeneity by succeeding if \emph{any} sub-goal closes,
transfers far more broadly.

\subsection{Vacuous compilation: a false positive mode}

Parts~8--9 identified a subtle failure: tactics that \emph{compile}
but whose distinctive component does zero work.
\texttt{exact\_mod\_cast} and \texttt{rw\_mod\_cast} compiled in
Representation Theory (Part~9, items~88--89) because the
\texttt{mod\_cast} preprocessor found no coercion lemmas to
apply---so it silently passed through, and the plain
\texttt{exact} or \texttt{rw} did all the work.  These are
\emph{not} successful transfers: the distinctive component (cast
normalization) contributed nothing.  The campaign criteria correctly
classify them as failures.

\subsection{Degenerate attractors}

Small target pools produce degenerate attractors: a few proof
states with high entity counts that score well against all source
queries, regardless of mathematical content.  In Model Theory
(Part~6), the \texttt{instSubsingleton} proof (3~entities) attracted
6 of 10~schemas.  In Condensed Mathematics (Part~8), three target
theorems attracted all 10~schemas.  Degenerate attractors inflate
QAP scores without reflecting genuine structural affinity, and
the resulting transfers often succeed trivially or fail
uninformatively.

%======================================================================
\section{Improving the General Theorem Prover}\label{sec:prover}
%======================================================================

The Yanasse campaign was a stress test of the Deep Vision matching
engine applied to formal mathematics.  Its findings directly
inform how the general theorem prover
(\texttt{proving\_loop.py}~\cite{linhares_deep_vision}) should
prioritize analogy sources, interpret QAP scores, and select tactics.

\subsection{Which areas should the prover look to for analogues?}

\begin{finding}[Source area prioritization]%
\label{find:source_priority}
When the prover encounters a proof state, it should prioritize
analogy matches from source areas whose tactic profiles complement
the target's needs:
\begin{enumerate}
\item \textbf{Probability} is the strongest general-purpose source
  (46\% across 7~targets).  Its structural combinators transfer
  broadly and operate on logical structure that every area shares.
  Probability should be the \emph{default} high-priority source
  for any target area.

\item \textbf{Number Theory} is a specialist source (10\% overall,
  but 30\% against Dynamics).  It should be prioritized \emph{only}
  when the target proof state carries algebraic type-class instances
  (\texttt{CommSemiring}, \texttt{Field},
  \texttt{LinearOrderedCommSemiring}).  A type-class pre-check can
  filter this efficiently.

\item \textbf{Analysis} (untested but predicted): the campaign's
  two-factor model predicts that Analysis---with its mix of
  structural combinators (\texttt{filter\_upwards},
  \texttt{continuity}) and computational tactics
  (\texttt{norm\_num}, \texttt{nlinarith})---would be a strong
  general-purpose source, possibly rivaling Probability.  It
  should be the next area tested.

\item \textbf{Algebra} (untested but predicted): with its rich
  decision procedures (\texttt{ring}, \texttt{group},
  \texttt{field\_simp}), Algebra would likely mirror Number
  Theory's profile: narrow but genuine transfers to targets with
  algebraic structure.
\end{enumerate}
\end{finding}

\subsection{Which areas should be de-prioritized?}

\begin{finding}[Target area difficulty]%
\label{find:target_difficulty}
The prover should adjust its confidence in analogy-guided tactics
based on the target area's position on the difficulty spectrum:
\begin{enumerate}
\item \textbf{Information Theory, Model Theory} ($\geq$50\%):
  analogy matching is highly effective; the prover should trust
  analogy-suggested tactics and invest reasoning budget in
  adaptation.

\item \textbf{Set Theory, Field Theory, Dynamics} (35--40\%):
  analogy matching works but is not dominant; the prover should
  use analogy suggestions as \emph{first attempts} but fall back
  to standard tactics quickly.

\item \textbf{Representation Theory} (20\%): analogy works only
  from structural sources (Probability); algebraic sources
  (Number Theory) are completely ineffective.  The prover should
  restrict analogy sources to structural-combinator-heavy areas.

\item \textbf{Condensed Mathematics} (10\%): analogy matching is
  nearly useless.  The prover should de-prioritize analogy-guided
  tactics and invest in category-theory-specific automation
  (\texttt{aesop\_cat}, categorical \texttt{simp} lemma sets)
  instead.

\item \textbf{Category Theory as source}: should be \emph{excluded}
  from analogy source pools.  Its DSL-style tactics
  (\texttt{cat\_disch}, \texttt{aesop\_cat}) encode
  domain-specific diagrammatic reasoning, not transferable proof
  strategies.  They dominated the initial rankings and had to be
  filtered out.
\end{enumerate}
\end{finding}

\subsection{Which QAP scores are informative?}

The QAP (Quadratic Assignment Problem) matching score is the
primary signal the prover uses to rank analogy candidates.  The
campaign reveals that its predictive value is real but bounded.

\begin{finding}[QAP score interpretation]%
\label{find:qap}
\begin{enumerate}
\item \textbf{Moderate scores (8--20) are the most predictive
  range.}  The majority of successful transfers fall in this range.
  Items~90 (QAP=14.49, \texttt{ring}$\to$Dynamics) and~92
  (QAP=12.17, \texttt{ring\_nf}$\to$Dynamics) are representative:
  the matcher identified genuine structural similarity between proof
  states involving real-valued functions with algebraic equality
  goals.

\item \textbf{Very high scores ($>$100) often signal degenerate
  attractors or pool artifacts,} not genuine structural affinity.
  Items~64 (QAP=125.0) and~66 (QAP=118.75) in Information Theory
  succeeded, but items with similarly high scores in Model Theory
  converged on a trivial 3-entity target.  The prover should treat
  scores above 100 with caution and cross-check against entity
  count.

\item \textbf{Zero scores can still yield successful transfers}
  for meta-tactics.  Item~93 (\texttt{enter}, QAP=0.0) succeeded
  because the tactic's utility depends on expression nesting depth,
  not on the entity-relation structure that QAP measures.  The
  prover should not discard zero-QAP candidates when the schema
  is a known meta-tactic.

\item \textbf{QAP is blind to type-class compatibility.}  A high
  QAP score does not guarantee that the source tactic's type-class
  requirements are met in the target.  Items~70 (QAP=17.75,
  \texttt{ring}$\to$Condensed) and~80 (QAP=17.19,
  \texttt{ring}$\to$RepTheory) scored well but failed completely.
  \emph{The prover should add a type-class pre-filter before
  investing reasoning budget in adaptation.}
\end{enumerate}
\end{finding}

\subsection{How could QAP matching be improved for Lean/Mathlib?}

The campaign exposes four specific limitations of the current
matching engine that, if addressed, would improve the prover's
performance on Lean proof states:

\begin{enumerate}
\item \textbf{Type-class-aware scoring.}
  The 14~relation types extracted from proof states
  (rewrite, fit/apply, head-match, structure, etc.)
  capture \emph{logical} relationships between entities but not
  \emph{algebraic} relationships between types.  Adding a relation
  type that encodes type-class instances
  (e.g., ``this type carries \texttt{CommSemiring}'') would allow
  the matcher to distinguish proof states where algebraic tactics
  can execute from those where they cannot---directly addressing
  the paradigm boundary.

\item \textbf{Entity-count normalization.}
  The current QAP score is unnormalized: larger proof states
  (more entities, more relations) produce higher raw scores.  This
  creates a bias toward matching large source states with large
  target states, and causes small but structurally rich targets to
  be overlooked.  Normalizing by $\sqrt{n_1 \cdot n_2}$ (geometric
  mean of entity counts) would make scores comparable across
  different-sized proof states.

\item \textbf{Nesting-depth features for meta-tactics.}
  Meta-tactics like \texttt{enter}, \texttt{conv}, and
  \texttt{show} operate on expression \emph{shape} rather than
  mathematical \emph{content}.  The current relation extractor does
  not capture nesting depth, application depth, or binder structure.
  Adding expression-shape features---even simple ones like
  ``maximum binder depth'' or ``number of nested function
  applications''---would let the matcher predict meta-tactic
  applicability.

\item \textbf{Area-aware weighting.}
  The campaign demonstrates that the same relation type
  (e.g., head-match) has different predictive value depending on
  the source-target pair.  Head-match is highly informative for
  Probability$\to$Information Theory (where shared measure-theoretic
  heads like \texttt{MeasureTheory.Measure} signal genuine affinity)
  but uninformative for Number Theory$\to$Condensed (where the heads
  are categorically different).  Learning area-specific relation
  weights from the campaign's 100~data points could improve matching
  precision.
\end{enumerate}

\subsection{Which tactics should be prioritized?}

\begin{finding}[Tactic prioritization for the prover]%
\label{find:tactic_priority}
Based on the 100-attempt campaign, the prover should adjust its
tactic exploration order as follows:

\paragraph{High priority (transfer broadly, succeed reliably):}
\begin{itemize}
\item \texttt{ext1}, \texttt{by\_cases}, \texttt{any\_goals}:
  succeeded in 5+ of 7~Probability targets.  These are the
  workhorses of cross-area transfer.
\item \texttt{congr with}, \texttt{push}: succeeded in 3--4
  targets.  Reliable when the target has equality or membership
  goals.
\item \texttt{ring}, \texttt{ring\_nf}: succeeded only against
  targets with \texttt{CommSemiring}, but when applicable, they
  close goals that no other tactic can.  Should be tried whenever
  the type-class pre-filter passes.
\end{itemize}

\paragraph{Medium priority (succeed in specific contexts):}
\begin{itemize}
\item \texttt{all\_goals}: succeeded in 3~targets (Dynamics,
  FieldTheory, InfoTheory) but fails on heterogeneous sub-goal
  sets.  The prover should check sub-goal homogeneity before
  trying.
\item \texttt{measurability}: succeeded only where
  $\sigma$-algebra structure exists (InfoTheory horizontal
  substitution for \texttt{fun\_prop}).
\item \texttt{enter}: domain-independent but requires expression
  nesting.  Check nesting depth, not QAP score.
\item \texttt{split\_ifs}: succeeded once (InfoTheory) on a
  piecewise-defined function.  Should be tried when the target
  contains \texttt{if-then-else} expressions.
\end{itemize}

\paragraph{Low priority (rarely transfer):}
\begin{itemize}
\item \texttt{fun\_prop}: 1~success out of 10~attempts (only
  where measurability goals were already present).  Domain-gated
  and highly specialized.
\item \texttt{lift}: 0~successes in 10~attempts.  Requires
  \texttt{CanLift} instances that are essentially unique to
  Probability's type ecosystem (\texttt{ENNReal}$\to$\texttt{NNReal},
  \texttt{WithTop}$\to$base type).
\item \texttt{congrm}: 1~success (InfoTheory).  Requires matched
  outer-operator patterns between goal sides, which is rare outside
  measure theory.
\item \texttt{decide}: 0~successes.  Requires ground propositions
  with no free variables---extremely rare in abstract mathematics.
\item \texttt{exact\_mod\_cast}, \texttt{rw\_mod\_cast}: 0~genuine
  successes (compiled twice but the cast component did zero work).
  Require numeric coercion chains absent outside Number Theory and
  Analysis.
\item \texttt{norm\_num}: 0~successes.  Requires numeric literal
  content in goal.
\item \texttt{linarith}: 0~successes outside its home areas.
  Requires linear-ordered commutative semiring, which is absent
  from most target types.
\item \texttt{field}: 0~successes.  Requires \texttt{Field}
  instance, absent from all tested targets.
\end{itemize}
\end{finding}

\subsection{Which tactics should be de-prioritized?}

The campaign identifies a clear hierarchy.  Tactics that \emph{never}
succeeded in cross-area transfer---\texttt{lift}, \texttt{decide},
\texttt{norm\_num}, \texttt{linarith}, \texttt{field},
\texttt{exact\_mod\_cast}, \texttt{rw\_mod\_cast}---should be
de-prioritized in the analogy-guided pipeline.  This does not mean
they are useless: \texttt{linarith} and \texttt{norm\_num} are
powerful \emph{within} their home areas.  But the analogy engine's
job is to find \emph{cross-area} connections, and these tactics are
too domain-specific to serve as transfer vehicles.

The prover already includes these tactics in its standard-tactics
fallback (step~7 of the pipeline).  The campaign confirms that this
is the right placement: they should be tried after domain-independent
strategies are exhausted, not promoted by analogy matching.

\subsection{Transfer maps from the paradigm boundary}

The metatheoretic results of Section~\ref{sec:metatheory} (Layers
1--3 and the \texttt{ConstantDependent} classification) are not
merely retrospective explanations of the campaign's failures.  They
provide a \emph{formal heuristic} that the prover can use
prospectively, before attempting any transfer.

By classifying each tactic's type-class prerequisites as either
\emph{constant-dependent} (requiring an element of the carrier)
or \emph{constant-free} (requiring only structural operations),
the prover can construct a \emph{transfer map}: a pre-computed
table indicating, for each (tactic class, target type) pair,
whether transfer is logically possible or provably impossible.
When the target type is empty or lacks the required algebraic
instances, constant-dependent tactics are pruned before the
expensive QAP matching and adaptation stages, saving the
reasoning budget for tactics that can actually succeed.

More concretely, when the prover encounters a new proof state, the
transfer map answers two questions immediately:

\begin{enumerate}
\item \textbf{When to stop trying to generalize an algebraic proof.}
  If the target type has no \texttt{CommSemiring} instance and the
  best library matches all use \texttt{ring} or \texttt{linarith},
  the transfer map says: these tactics are provably inapplicable.
  Do not invest reasoning budget in semantic adaptation.  Fall back
  to structural tactics or domain-specific lemmas.

\item \textbf{When to extract the structural skeleton.}
  If an algebraic proof cannot transfer whole, the transfer map
  identifies which \emph{sub-steps} are structural (and therefore
  transferable) and which are algebraic (and therefore blocked).
  The prover can then transfer the structural skeleton---the
  case splits, extensionality reductions, and congruence steps---while
  replacing the algebraic sub-steps with target-appropriate tactics.
  This is precisely the (head, modifier) decomposition of
  Section~\ref{sec:transfers}: the modifier (structural skeleton)
  crosses the boundary; the head (algebraic decision procedure)
  does not.
\end{enumerate}

%======================================================================
\section{The Two-Factor Model}\label{sec:model}
%======================================================================

The campaign's results are well described by a two-factor model:

\[
P(\text{transfer success}) \;\approx\;
  f(\text{source tactic profile}) \;\times\;
  g(\text{target type-class landscape})
\]

\noindent where:

\begin{itemize}
\item $f$ measures how \emph{domain-general} the source tactic is.
  Structural combinators ($f \approx 0.6$--$0.7$) transfer broadly;
  algebraic decision procedures ($f \approx 0.1$) transfer narrowly.

\item $g$ measures how \emph{type-class-accessible} the target is.
  Information Theory ($g \approx 0.9$) shares Probability's
  measure-theoretic infrastructure; Condensed Mathematics
  ($g \approx 0.1$) has an entirely categorical type-class landscape.
\end{itemize}

The model explains the key patterns:

\begin{enumerate}
\item \textbf{Probability$\to$InfoTheory (90\%):}
  High $f$ (structural combinators) $\times$ high $g$ (shared
  measure-theoretic types) = very high transfer rate.

\item \textbf{NumTh$\to$Condensed (0\%):}
  Low $f$ (algebraic procedures) $\times$ low $g$ (no algebraic
  instances) = complete failure.

\item \textbf{NumTh$\to$Dynamics (30\%):}
  Low $f$ $\times$ moderate $g$ (partial $\mathbb{R}$-arithmetic
  sub-layers) = partial success, concentrated in specific sub-areas.

\item \textbf{Prob$\to$Condensed (20\%):}
  High $f$ $\times$ low $g$ = limited success, only from the most
  domain-agnostic tactics (\texttt{ext1}, \texttt{congr with}).
\end{enumerate}

The model also predicts untested pairs.  For instance,
Probability$\to$Analysis should achieve a very high rate (structural
combinators + shared measure-theoretic types + Analysis has rich
algebraic structure).  Number Theory$\to$Analysis should achieve a
moderate rate (algebraic procedures + Analysis has
\texttt{CommSemiring} everywhere).

\subsection{Complementary sources}

The two-factor model implies that different sources access
different parts of the same target.  This was confirmed empirically
for Dynamics: Probability transferred structural combinators to the
measure-theoretic and extensionality layers (4~new theorems);
Number Theory transferred algebraic procedures to the arithmetic
core (3~new theorems).  The two sets of new theorems are entirely
non-overlapping.

A multi-source campaign is therefore strictly more powerful than any
single-source campaign.  The prover should maintain analogy
libraries from multiple source areas and query them in parallel,
weighted by the two-factor model's prediction for the current
target.

%======================================================================
\section{The Structural Map of Mathlib Transferability}%
\label{sec:map}
%======================================================================

The campaign's data, combined with the transferability-weighted
model of Section~\ref{sec:model}, allows us to compute the full
transfer potential between every pair of Mathlib's 23~mathematical
areas (excluding \texttt{CategoryTheory} as source and
\texttt{Tactic}/\texttt{Control}/\texttt{Logic} as targets;
see Part~1~\cite{yanasse_part1} for filtering rationale).  We present two matrices:
the \emph{naive} model (pair potential computed from raw $z$-score
gaps, used to select the initial campaign's 10~pairs) and the
\emph{transferability-weighted} model (pair potential weighted by
empirical transfer rates per tactic class, derived from the
campaign's 100~data points).

Table~\ref{tab:legend} provides the area index used in both
matrices.  Tables~\ref{tab:naive_matrix} and
\ref{tab:weighted_matrix} present the $23 \times 23$ transfer
potential matrices.  Each cell shows the pair potential (sum of
top-10 schema transfer gaps, weighted by $\log(1 + \text{source
count})$ and, in the weighted model, multiplied by the tactic
class's empirical transfer probability).  Higher values indicate
greater theoretical transfer opportunity.  The diagonal is
excluded (self-transfer is trivial).

\subsection{Naive vs.\ weighted rankings}

The reranking is dramatic.  In the naive model, Number Theory
pairs dominate the top~20 (high $z$-score gaps from algebraic
decision procedures).  In the weighted model, those pairs drop
sharply because the campaign showed algebraic procedures transfer
at ${\sim}$10\% vs.\ ${\sim}$60\% for structural combinators.
Probability remains the top source in both models, but the
weighted model promotes new sources---Measure Theory, Algebraic
Topology, Representation Theory---whose structural combinators
were buried under Number Theory's inflated gaps.

The full matrices and the underlying data
(\texttt{area\_pair\_potential.parquet},
\texttt{transfer\_candidates.parquet}, \texttt{schemas.parquet})
are available in the project's GitHub repository.

\begin{table}[ht]
\centering\small
\caption{Area index for the transfer potential matrices.}
\label{tab:legend}
\begin{tabular}{@{}rll@{}}
\toprule
\# & Area & Mathlib namespace \\
\midrule
1 & Algebra & \texttt{Mathlib.Algebra} \\
2 & Algebraic Geometry & \texttt{Mathlib.AlgebraicGeometry} \\
3 & Algebraic Topology & \texttt{Mathlib.AlgebraicTopology} \\
4 & Analysis & \texttt{Mathlib.Analysis} \\
5 & Combinatorics & \texttt{Mathlib.Combinatorics} \\
6 & Computability & \texttt{Mathlib.Computability} \\
7 & Condensed & \texttt{Mathlib.Condensed} \\
8 & Data & \texttt{Mathlib.Data} \\
9 & Dynamics & \texttt{Mathlib.Dynamics} \\
10 & Field Theory & \texttt{Mathlib.FieldTheory} \\
11 & Geometry & \texttt{Mathlib.Geometry} \\
12 & Group Theory & \texttt{Mathlib.GroupTheory} \\
13 & Information Theory & \texttt{Mathlib.InformationTheory} \\
14 & Linear Algebra & \texttt{Mathlib.LinearAlgebra} \\
15 & Measure Theory & \texttt{Mathlib.MeasureTheory} \\
16 & Model Theory & \texttt{Mathlib.ModelTheory} \\
17 & Number Theory & \texttt{Mathlib.NumberTheory} \\
18 & Order & \texttt{Mathlib.Order} \\
19 & Probability & \texttt{Mathlib.Probability} \\
20 & Representation Theory & \texttt{Mathlib.RepresentationTheory} \\
21 & Ring Theory & \texttt{Mathlib.RingTheory} \\
22 & Set Theory & \texttt{Mathlib.SetTheory} \\
23 & Topology & \texttt{Mathlib.Topology} \\
\bottomrule
\end{tabular}
\end{table}

\begin{landscape}
\begin{table}[p]
\centering
\caption{Naive transfer potential matrix (raw $z$-score gaps, no transferability weighting).
  Row = source area, column = target area.  The initial campaign selected the
  10~highest-potential pairs from this matrix.}
\label{tab:naive_matrix}
\setlength{\tabcolsep}{2.5pt}
{\fontsize{5.5}{7}\selectfont
\begin{tabular}{@{}r|rrrrrrrrrrrrrrrrrrrrrrr@{}}
\toprule
& 1 & 2 & 3 & 4 & 5 & 6 & 7 & 8 & 9 & 10 & 11 & 12 & 13 & 14 & 15 & 16 & 17 & 18 & 19 & 20 & 21 & 22 & 23 \\
\midrule
1 & -- & 33 & 33 & 23 & 33 & 33 & 49 & 25 & 49 & 49 & 49 & 49 & 49 & 23 & 46 & 49 & 23 & 23 & 33 & 49 & 7 & 49 & 11 \\
2 & 74 & -- & 152 & 86 & 134 & 166 & 175 & 119 & 162 & 122 & 99 & 142 & 175 & 121 & 142 & 170 & 98 & 140 & 150 & 167 & 56 & 151 & 118 \\
3 & 34 & 70 & -- & 7 & 61 & 103 & 141 & 38 & 149 & 82 & 42 & 90 & 137 & 31 & 84 & 116 & 42 & 58 & 96 & 79 & 31 & 107 & 33 \\
4 & 79 & 120 & 120 & -- & 74 & 114 & 129 & 118 & 129 & 116 & 80 & 106 & 129 & 118 & 84 & 129 & 60 & 114 & 99 & 129 & 74 & 129 & 28 \\
5 & 49 & 92 & 76 & 43 & -- & 89 & 127 & 45 & 116 & 98 & 97 & 72 & 127 & 82 & 83 & 115 & 81 & 68 & 92 & 127 & 82 & 104 & 72 \\
6 & 47 & 135 & 138 & 59 & 90 & -- & 165 & 58 & 143 & 134 & 116 & 98 & 152 & 111 & 89 & 143 & 77 & 118 & 127 & 154 & 73 & 143 & 60 \\
7 & 8 & 26 & 43 & 13 & 20 & 64 & -- & 23 & 64 & 20 & 47 & 20 & 101 & 8 & 21 & 55 & 28 & 20 & 44 & 71 & 8 & 55 & \\
8 & 55 & 120 & 138 & 78 & 65 & 86 & 165 & -- & 145 & 90 & 118 & 91 & 165 & 84 & 75 & 129 & 93 & 75 & 119 & 145 & 50 & 93 & 87 \\
9 & 28 & 62 & 74 & 5 & 47 & 82 & 89 & 51 & -- & 55 & 68 & 60 & 87 & 56 & 21 & 83 & 35 & 42 & 60 & 90 & 14 & 60 & 26 \\
10 & 41 & 81 & 120 & 41 & 82 & 92 & 163 & 78 & 137 & -- & 77 & 75 & 140 & 49 & 75 & 110 & 46 & 98 & 110 & 118 & 24 & 100 & 62 \\
11 & 59 & 86 & 83 & 11 & 84 & 91 & 117 & 64 & 91 & 86 & -- & 82 & 117 & 48 & 66 & 101 & 45 & 75 & 83 & 113 & 65 & 87 & 55 \\
12 & 33 & 80 & 91 & 61 & 64 & 87 & 91 & 54 & 91 & 88 & 67 & -- & 91 & 69 & 81 & 89 & 66 & 85 & 84 & 91 & 36 & 85 & 61 \\
13 & 19 & 23 & 46 & & 23 & 36 & 61 & 41 & 54 & 48 & 3 & 36 & -- & 41 & 19 & 56 & & 65 & 15 & 68 & 23 & 52 & 20 \\
14 & 20 & 90 & 107 & 53 & 80 & 96 & 107 & 25 & 107 & 78 & 62 & 57 & 107 & -- & 68 & 107 & 60 & 102 & 85 & 107 & 31 & 107 & 55 \\
15 & 102 & 137 & 147 & 7 & 88 & 70 & 147 & 147 & 147 & 147 & 51 & 147 & 125 & 110 & -- & 147 & 76 & 64 & 27 & 147 & 68 & 137 & 66 \\
16 & 25 & 52 & 76 & 25 & 32 & 64 & 82 & 29 & 82 & 55 & 52 & 57 & 82 & 35 & 29 & -- & 41 & 60 & 61 & 82 & 20 & 64 & 29 \\
17 & 76 & 128 & 144 & 33 & 112 & 135 & 189 & 85 & 189 & 147 & 121 & 123 & 173 & 99 & 70 & 183 & -- & 177 & 110 & 189 & 48 & 176 & 92 \\
18 & 29 & 84 & 84 & 34 & 60 & 81 & 94 & 41 & 94 & 67 & 67 & 60 & 94 & 58 & 59 & 90 & 60 & -- & 76 & 94 & 48 & 84 & 39 \\
19 & 70 & 132 & 187 & 3 & 116 & 170 & 208 & 144 & 204 & 207 & 129 & 171 & 190 & 162 & 31 & 204 & 86 & 137 & -- & 219 & 101 & 209 & 72 \\
20 & 17 & 37 & 59 & 9 & 27 & 77 & 129 & 43 & 105 & 49 & 24 & 40 & 133 & 16 & 20 & 100 & 28 & 31 & 79 & -- & 6 & 106 & 30 \\
21 & 18 & 55 & 78 & 12 & 49 & 78 & 93 & 22 & 93 & 60 & 32 & 39 & 93 & 40 & 56 & 78 & 25 & 70 & 70 & 93 & -- & 70 & 27 \\
22 & 35 & 82 & 117 & 38 & 72 & 101 & 149 & 17 & 113 & 100 & 82 & 82 & 128 & 81 & 78 & 132 & 79 & 66 & 108 & 160 & 56 & -- & 46 \\
23 & 52 & 104 & 145 & 27 & 96 & 139 & 146 & 72 & 146 & 104 & 92 & 110 & 146 & 92 & 75 & 104 & 102 & 87 & 114 & 146 & 82 & 132 & -- \\
\bottomrule
\end{tabular}
}
\end{table}

\begin{table}[p]
\centering
\caption{Transferability-weighted potential matrix.  Same structure as
  Table~\ref{tab:naive_matrix}, but each schema's gap is multiplied by
  its tactic class's empirical transfer probability (from the campaign's
  100~data points).  This model should be used for selecting future
  transfer pairs.}
\label{tab:weighted_matrix}
\setlength{\tabcolsep}{2.5pt}
{\fontsize{5.5}{7}\selectfont
\begin{tabular}{@{}r|rrrrrrrrrrrrrrrrrrrrrrr@{}}
\toprule
& 1 & 2 & 3 & 4 & 5 & 6 & 7 & 8 & 9 & 10 & 11 & 12 & 13 & 14 & 15 & 16 & 17 & 18 & 19 & 20 & 21 & 22 & 23 \\
\midrule
1 & -- & 9 & 9 & 7 & 9 & 9 & 13 & 7 & 13 & 13 & 13 & 13 & 13 & 7 & 12 & 13 & 7 & 7 & 9 & 13 & 2 & 13 & 3 \\
2 & 23 & -- & 40 & 25 & 35 & 56 & 53 & 33 & 55 & 40 & 27 & 39 & 53 & 34 & 39 & 51 & 24 & 38 & 46 & 44 & 18 & 47 & 34 \\
3 & 9 & 20 & -- & 4 & 17 & 36 & 45 & 10 & 47 & 31 & 13 & 25 & 44 & 10 & 30 & 31 & 13 & 16 & 26 & 22 & 7 & 37 & 9 \\
4 & 21 & 41 & 41 & -- & 20 & 35 & 41 & 38 & 41 & 38 & 28 & 36 & 41 & 39 & 32 & 41 & 24 & 35 & 32 & 41 & 28 & 41 & 15 \\
5 & 15 & 30 & 26 & 12 & -- & 32 & 40 & 13 & 35 & 31 & 27 & 21 & 40 & 25 & 28 & 37 & 26 & 20 & 30 & 40 & 26 & 33 & 22 \\
6 & 12 & 33 & 38 & 15 & 23 & -- & 45 & 15 & 40 & 36 & 34 & 27 & 43 & 31 & 25 & 35 & 21 & 27 & 37 & 41 & 18 & 36 & 15 \\
7 & & 5 & 16 & 4 & 9 & 24 & -- & 5 & 20 & 9 & 15 & 9 & 34 & & 7 & 18 & 8 & 9 & 10 & 22 & & 17 & \\
8 & 14 & 32 & 34 & 22 & 19 & 24 & 43 & -- & 38 & 25 & 32 & 25 & 43 & 23 & 21 & 35 & 26 & 21 & 32 & 38 & 12 & 26 & 22 \\
9 & 5 & 15 & 18 & 2 & 10 & 23 & 24 & 11 & -- & 12 & 17 & 14 & 24 & 12 & 3 & 23 & 9 & 10 & 14 & 25 & 3 & 16 & 5 \\
10 & 10 & 20 & 30 & 10 & 20 & 23 & 41 & 20 & 34 & -- & 19 & 19 & 35 & 13 & 19 & 28 & 12 & 24 & 28 & 30 & 6 & 25 & 15 \\
11 & 13 & 19 & 19 & 3 & 16 & 21 & 28 & 15 & 21 & 19 & -- & 16 & 28 & 10 & 12 & 23 & 12 & 14 & 18 & 27 & 14 & 20 & 8 \\
12 & 8 & 21 & 24 & 17 & 16 & 23 & 24 & 13 & 24 & 23 & 17 & -- & 24 & 19 & 22 & 24 & 17 & 21 & 22 & 24 & 9 & 23 & 17 \\
13 & 2 & 3 & 8 & & 2 & 6 & 18 & 4 & 17 & 3 & & 4 & -- & 3 & 2 & 5 & & 8 & 1 & 8 & 2 & 4 & 5 \\
14 & 5 & 27 & 32 & 13 & 25 & 29 & 32 & 8 & 32 & 21 & 17 & 16 & 32 & -- & 18 & 32 & 20 & 27 & 23 & 32 & 9 & 32 & 15 \\
15 & 36 & 47 & 50 & 2 & 32 & 27 & 50 & 50 & 50 & 50 & 11 & 50 & 43 & 30 & -- & 50 & 29 & 25 & 8 & 50 & 18 & 47 & 26 \\
16 & 7 & 15 & 20 & 6 & 8 & 17 & 21 & 8 & 21 & 15 & 13 & 15 & 21 & 9 & 8 & -- & 11 & 16 & 16 & 21 & 5 & 17 & 8 \\
17 & 16 & 26 & 32 & 11 & 22 & 31 & 33 & 17 & 32 & 31 & 25 & 25 & 32 & 22 & 19 & 29 & -- & 28 & 23 & 32 & 10 & 32 & 20 \\
18 & 11 & 35 & 36 & 12 & 23 & 29 & 38 & 12 & 38 & 24 & 24 & 20 & 38 & 20 & 19 & 35 & 24 & -- & 25 & 38 & 16 & 35 & 10 \\
19 & 29 & 35 & 68 & & 37 & 68 & 92 & 37 & 90 & 78 & 43 & 57 & 87 & 54 & 10 & 79 & 36 & 35 & -- & 92 & 21 & 85 & 29 \\
20 & 4 & 9 & 15 & 2 & 7 & 30 & 38 & 11 & 43 & 18 & 6 & 10 & 54 & 4 & 5 & 31 & 7 & 8 & 25 & -- & 1 & 48 & 7 \\
21 & 4 & 13 & 19 & 4 & 11 & 19 & 23 & 4 & 23 & 15 & 8 & 9 & 23 & 10 & 14 & 19 & 8 & 17 & 19 & 23 & -- & 17 & 5 \\
22 & 8 & 19 & 32 & 9 & 16 & 25 & 41 & 4 & 30 & 29 & 24 & 19 & 36 & 21 & 19 & 33 & 20 & 16 & 30 & 43 & 14 & -- & 11 \\
23 & 14 & 28 & 38 & 8 & 24 & 36 & 38 & 20 & 38 & 28 & 26 & 30 & 38 & 26 & 20 & 28 & 29 & 24 & 30 & 38 & 19 & 35 & -- \\
\bottomrule
\end{tabular}
}
\end{table}
\end{landscape}

%======================================================================
\section{Toward a Metatheoretic Impossibility Result}%
\label{sec:metatheory}
%======================================================================

The paradigm boundary identified in Sections~\ref{sec:barriers}
and~\ref{sec:model} is an empirical finding: 20~attempts by Number
Theory's algebraic tactics against Condensed Mathematics and
Representation Theory produced 0~successes.  But the uniformity of
the failure mode---every attempt traces to a missing type-class
instance---suggests that the barrier is not an engineering limitation
but a \emph{logical necessity}.  Can this be proved?

\subsection{The claim}

\begin{quote}
No tactic whose proof terms are built from the axioms of
\texttt{CommSemiring} can close a goal whose type carries no
\texttt{CommSemiring} instance.
\end{quote}

\noindent This is a proof-theoretic statement: the derivable
equalities of the algebraic fragment form a restricted class, and
morphism-level equalities are provably outside it.

\subsection{Three layers of formalization}

\paragraph{Layer~1: Non-existence of instances.}
Prove in Lean that specific types lack specific instances:
\begin{lstlisting}
-- Morphisms in an abstract category carry no CommSemiring
theorem no_CommSemiring_hom {C : Type*} [Category C]
    (X Y : C) (hne : X (*$\neq$*) Y) :
    (*$\neg$*) CommSemiring (X (*$\longrightarrow$*) Y) := ...
\end{lstlisting}
Combined with the documented precondition that \texttt{ring}
requires \texttt{CommSemiring}, this yields: ``\texttt{ring} cannot
close goals of type $f = g$ where $f, g : X \longrightarrow Y$ in
an abstract category.''  This is standard Lean mathematics---no
metaprogramming reflection required.

\paragraph{Layer~2: Proof-term fragment analysis.}
Formalize \emph{what proof terms \texttt{ring} can produce}.
Mathlib's \texttt{Mathlib.Tactic.Ring} constructs proof terms via a
ring normal form: compositions of \texttt{add\_comm},
\texttt{mul\_assoc}, \texttt{add\_zero}, \texttt{mul\_one},
\texttt{distrib}, and their consequences.  Define this fragment as
an inductive type:
\begin{lstlisting}
inductive RingDerivable [CommSemiring R] :
    R (*$\to$*) R (*$\to$*) Prop where
  | refl : RingDerivable a a
  | add_comm : RingDerivable (a + b) (b + a)
  | mul_assoc : RingDerivable (a * b * c) (a * (b * c))
  | trans : RingDerivable a b (*$\to$*) RingDerivable b c
      (*$\to$*) RingDerivable a c
  ...
\end{lstlisting}
The impossibility follows from the type signature: \texttt{RingDerivable}
has \texttt{[CommSemiring R]} as a parameter.  Without the instance,
the inductive type \emph{cannot be instantiated}---in dependent type
theory, the paradigm boundary is enforced by the type system itself.
The theorem is: ``the \texttt{ring} tactic produces a proof by
constructing a \texttt{RingDerivable} witness and extracting
$a = b$ via soundness; constructing the witness requires
\texttt{CommSemiring}; therefore, if the goal type has no such
instance, \texttt{ring} cannot produce a proof.''

\paragraph{Layer~3: General paradigm boundary.}
The most ambitious version: for \emph{any} tactic whose proof terms
live in an axiomatic fragment~$F$, and any goal type~$T$ that doesn't
model $F$'s axioms, the tactic cannot close the goal.  This is
analogous to separation results in proof complexity (Resolution vs.\
Cutting Planes vs.\ Frege systems).  The algebraic fragment
(\texttt{CommSemiring} axioms) and the categorical fragment
(composition, identity, functoriality axioms) generate
\emph{non-overlapping} sets of derivable equalities---except where a
type happens to model both (e.g., endomorphism rings).

\subsection{Why this matters}

A formal impossibility result would upgrade the paradigm boundary
from an empirical observation ($n = 20$, $p = 0$) to a mathematical
theorem.  It would also establish that the tactic taxonomy of
Section~\ref{sec:transfers}---structural combinators vs.\ algebraic
procedures vs.\ domain-specific lemmas---is not merely a practical
classification but reflects genuine proof-theoretic distinctions
between the axiomatic fragments that different tactic classes exploit.

The deeper question is whether the barrier is \emph{fundamentally}
about type-class instances or about something more structural.
Morphisms compose but do not add or multiply; composition is
non-commutative and non-distributive; there is no zero morphism in a
general category.  The algebraic axioms that ring tactics exploit are
\emph{structurally incompatible} with the categorical axioms that
morphism-level proofs use.  Formalizing this incompatibility---not
just for \texttt{ring} vs.\ morphisms, but for algebraic vs.\
categorical proof theory in general---would be a contribution to
the foundations of formalized mathematics.

\subsection{Layer~1: Formal proof}

We proved Layer~1 in Lean~4 with Mathlib.  The core theorem:

\begin{theorem}[No commutative semiring on empty types]%
\label{thm:paradigm}
Let $\alpha$ be a type with no inhabitants.  Then $\alpha$ admits no
commutative semiring structure.
\end{theorem}

\begin{proof}
Every commutative semiring has a zero element.  If $\alpha$ carried
a commutative semiring structure, it would therefore contain at
least one element (namely~$0$).  But $\alpha$ is empty.
Contradiction.
\end{proof}

\noindent The same argument generalizes to any algebraic type class
that provides an element: \texttt{Field}, \texttt{LinearOrdered\-CommSemiring},
\texttt{Ring}, or any class extending \texttt{Zero} or \texttt{One}.
The paradigm boundary is not specific to \texttt{CommSemiring}---it
applies to the entire algebraic hierarchy.

\subsection{The proving loop's search}

Before writing the proof by hand, we ran the Deep Vision proving
loop (\texttt{proving\_loop.py}) on the theorem statement---the same
relational analogy engine that discovered the paradigm boundary
empirically.  The system matched the proof state against 45{,}096
library entries via GPU-accelerated QAP on MPS, then explored a
proof search tree.  The key steps (abbreviated):

\medskip
\begin{center}\small
\begin{tabular}{@{}cp{4.5cm}p{5.5cm}@{}}
\toprule
D & Tactic (verbatim from trace) & Effect \\
\midrule
0 & \texttt{simp}
  & Unfolds the goal; progress, 5~entities remain \\[2pt]
1 & \texttt{intro h\_comm}
  & Assumes CommSemiring (contradiction setup) \\[2pt]
2 & \texttt{have := IsEmpty.elim h 0}
  & Tries to extract a zero element \\[2pt]
3 & \texttt{have := zero\_ne\_one; exact this this}
  & Tries the $0 \neq 1$ route \\[2pt]
4 & \texttt{rw [this] at *}
  & Cycle detected, backtrack \\[2pt]
5 & \texttt{haveI : Nontrivial := @CommSemiring.toNontrivial ...}
  & CommSemiring implies Nontrivial \\[2pt]
5 & \texttt{haveI : Subsingleton := IsEmpty.subsingleton}
  & IsEmpty implies Subsingleton \\[2pt]
& \multicolumn{2}{@{}l}{\emph{Max depth 6 reached; backtracking.}} \\
\bottomrule
\end{tabular}
\end{center}

\medskip
\noindent The Claude planning module also proposed the correct strategy
at depth~0: \texttt{intro inst; exact IsEmpty.false\_of\_nonempty h
$\langle$(0\,:\,$\alpha$)$\rangle$}.  The system found the right
pieces---\texttt{intro} the instance, extract a zero element, contradict
\texttt{IsEmpty}---but could not assemble them within the depth limit.

The critical obstacle was that after \texttt{intro inst}, the locally
introduced \texttt{inst} is not registered as a type-class instance by
Lean's synthesizer, so \texttt{(0\,:\,$\alpha$)} fails to elaborate.
The proof requires the explicit annotation
\texttt{@Zero.zero\,$\alpha$\,inst.toZero} to route through the
introduced hypothesis.  This is a gap in the proving loop's
adaptation layer: it can discover the mathematical strategy
(assume, extract, contradict) but cannot yet perform the
type-class threading needed to compile it.

\subsection{Verified Lean proof}

The following proof compiles with Lean~4.29.0 and Mathlib~4.29.0
(exit code~0, zero \texttt{sorry} declarations):

\begin{lstlisting}
import Mathlib

-- Theorem 1: An empty type cannot carry CommSemiring.
theorem no_CommSemiring_of_empty
    ((*$\alpha$*) : Type*) (h : IsEmpty (*$\alpha$*)) :
    CommSemiring (*$\alpha$*) (*$\to$*) False := by
  intro inst
  exact h.false (@Zero.zero (*$\alpha$*) inst.toZero)

-- Theorem 2: Same via (*$\neg$*) Nonempty.
theorem no_CommSemiring_of_not_nonempty
    ((*$\alpha$*) : Type*) (h : (*$\neg$*) Nonempty (*$\alpha$*)) :
    CommSemiring (*$\alpha$*) (*$\to$*) False := by
  intro inst
  exact h (*$\langle$*)@Zero.zero (*$\alpha$*) inst.toZero(*$\rangle$*)
\end{lstlisting}

\noindent The self-referential character of this result deserves
emphasis: \emph{the same relational analogy engine that discovered
the paradigm boundary empirically was then asked to prove that the
boundary is logically necessary}.  It found the right proof strategy
through structural matching against its library of 217{,}133 proof
states, but could not compile the final proof term due to a
type-class threading limitation.  The formal proof was completed by
hand, using the strategy the engine discovered.

\subsection{Layer~2: The ring proof fragment}

Layer~1 showed that empty types cannot carry \texttt{CommSemiring}.
Layer~2 goes further: it defines the \emph{equational theory} of
commutative semirings as an inductive type
\texttt{RingDerivable}---one constructor per axiom (commutativity,
associativity, distributivity, identity elements, congruence)---and
proves two results.

\begin{theorem}[Soundness of the ring fragment]%
\label{thm:ring_sound}
Every equality derivable from the commutative semiring axioms is
true.  That is, if $a =_{\mathrm{ring}} b$ is derivable, then
$a = b$.
\end{theorem}

\begin{proof}
By induction on the derivation.  Each constructor of
\texttt{RingDerivable} corresponds to a \texttt{CommSemiring} axiom
($a + b = b + a$, $a \cdot 1 = a$, etc.), and each axiom holds in
any commutative semiring.  Transitivity, symmetry, and congruence
preserve equality.
\end{proof}

\begin{theorem}[The ring fragment requires \texttt{CommSemiring}]%
\label{thm:ring_requires}
The type \texttt{RingDerivable} has \texttt{[CommSemiring\,R]} as an
instance parameter.  Without this instance, the type cannot be
instantiated: no ring-derivable equality can be \emph{stated}, let
alone proved.  Combined with Theorem~\ref{thm:paradigm}, this
yields: over an empty type, the ring proof fragment is vacuous.
\end{theorem}

\begin{proof}
In dependent type theory, an instance parameter is a hard
constraint enforced by the type checker.  Every constructor of
\texttt{RingDerivable} uses the operations $+$, $\times$, $0$, $1$
provided by \texttt{CommSemiring}.  Without the instance, these
operations are undefined and the constructors are ill-typed.
The type checker rejects any attempt to form a
\texttt{RingDerivable} term.
\end{proof}

\noindent The verified Lean formalization defines \texttt{RingDerivable}
with 17~constructors (reflexivity, symmetry, transitivity, and one
constructor per \texttt{CommSemiring} axiom plus congruence for $+$
and~$\times$), proves soundness by induction, and extends the
impossibility to \texttt{Field}, \texttt{Ring}, and \texttt{Semiring}
via the common root cause: all extend \texttt{Zero}, which provides
an element.

\begin{lstlisting}
-- The equational theory of commutative semirings (excerpt)
inductive RingDerivable {R : Type*} [CommSemiring R]
    : R (*$\to$*) R (*$\to$*) Prop where
  | rd_refl (a : R) : RingDerivable a a
  | rd_add_comm (a b : R) :
      RingDerivable (a + b) (b + a)
  | rd_mul_zero (a : R) :
      RingDerivable (a * 0) 0
  | rd_left_distrib (a b c : R) :
      RingDerivable (a * (b + c)) (a * b + a * c)
  ...  -- 17 constructors total

-- Soundness
theorem RingDerivable.sound {R : Type*}
    [CommSemiring R] {a b : R}
    (h : RingDerivable a b) : a = b := by
  induction h with
  | rd_add_comm a b => exact add_comm a b
  | rd_mul_zero a => exact mul_zero a
  ...  -- one case per constructor

-- The root cause: Zero provides an element
theorem no_Zero_of_empty ((*$\alpha$*) : Type*)
    (h : IsEmpty (*$\alpha$*)) (inst : Zero (*$\alpha$*)) : False :=
  h.false (@Zero.zero (*$\alpha$*) inst)

-- Every algebraic class reduces to this:
theorem no_Ring_of_empty ((*$\alpha$*) : Type*)
    (h : IsEmpty (*$\alpha$*)) (inst : Ring (*$\alpha$*)) : False :=
  no_Zero_of_empty (*$\alpha$*) h inst.toZero
\end{lstlisting}

\subsection{Layer~3: The general paradigm boundary}

The deepest formalization places the algebraic and structural
fragments side by side and proves they have fundamentally different
type-class requirements.

\begin{theorem}[The structural fragment is unrestricted]%
\label{thm:struct_unrestricted}
Define \texttt{StructDeriv} as the equational theory of pure logic:
reflexivity, symmetry, transitivity, and congruence under function
application.  This inductive type has \emph{no type-class parameters}.
It can be stated and proved over any type, including empty ones.
In particular, any two functions $f, g : \varnothing \to \beta$
are structurally equal (vacuously).
\end{theorem}

\begin{proof}
\texttt{StructDeriv} is parameterized only by a type
$\alpha$---no algebraic structure is required.  For functions out
of the empty type, the function-extensionality constructor
requires proving $f(x) = g(x)$ for all $x : \varnothing$, which
holds vacuously (there are no such~$x$).
\end{proof}

\begin{theorem}[The paradigm boundary is the constant-symbol interface]%
\label{thm:general_boundary}
A proof fragment whose signature includes a \emph{constant symbol}
(a nullary operation producing an element of the carrier---such as
$0$ or $1$) is incompatible with empty types.  A proof fragment
whose signature includes only \emph{function symbols} (operations
taking elements to elements, but not creating them) is compatible
with empty types.  The algebraic fragment is the former; the
structural fragment is the latter.  The paradigm boundary is
the interface between them.
\end{theorem}

\begin{proof}
A constant symbol $c$ of type $\alpha$ is an element of $\alpha$.
If $\alpha$ is empty, no such element exists.  Therefore any type
class providing a constant symbol (every class extending
\texttt{Zero} or \texttt{One}) contradicts emptiness.

Conversely, function symbols $f : \alpha^n \to \alpha$ with
$n \geq 1$ preserve but do not create elements: if $\alpha$ is
empty, there are no inputs, so no outputs need exist.  The
structural fragment's operations (function application, abstraction,
extensionality) are all of this kind.
\end{proof}

\noindent The verified Lean formalization defines both
\texttt{RingDeriv} (with \texttt{[CommSemiring\,R]}) and
\texttt{StructDeriv} (with no type-class parameters), proves
both sound, demonstrates \texttt{StructDeriv} over \texttt{Empty},
and shows that the entire algebraic hierarchy (\texttt{CommSemiring},
\texttt{Field}, \texttt{Ring}, \texttt{Semiring}) reduces to a
single root cause: \texttt{Zero} provides an element.

\begin{lstlisting}
-- The structural fragment: NO type-class parameters
inductive StructDeriv {(*$\alpha$*) : Type*}
    : (*$\alpha$*) (*$\to$*) (*$\alpha$*) (*$\to$*) Prop where
  | sd_refl (a : (*$\alpha$*)) : StructDeriv a a
  | sd_symm : StructDeriv a b (*$\to$*) StructDeriv b a
  | sd_trans : StructDeriv a b (*$\to$*) StructDeriv b c
      (*$\to$*) StructDeriv a c

-- Structural fragment WORKS on empty types:
theorem structural_vacuous ((*$\beta$*) : Type*)
    (f g : Empty (*$\to$*) (*$\beta$*)) : FunStructDeriv f g :=
  .fsd_funext fun x => Empty.elim x

-- Algebraic fragment BLOCKED on empty types:
theorem algebraic_blocked {(*$\alpha$*) : Type*}
    (hempty : IsEmpty (*$\alpha$*)) [CommSemiring (*$\alpha$*)]
    {a b : (*$\alpha$*)} (_h : RingDeriv a b) : False :=
  hempty.false a
\end{lstlisting}

\subsection{Machine-checkable boundary classification}

The three layers above prove the paradigm boundary for specific
type classes one at a time.  We unify them into a single
type-class hierarchy that makes the boundary \emph{machine-checkable}:
registering any new algebraic class automatically yields the
impossibility theorem.

\begin{definition}[\texttt{ConstantDependent}]%
\label{def:constant_dependent}
A type class $C$ is \emph{constant-dependent} if every instance of
$C$ on a type $\alpha$ provides at least one element of $\alpha$
(a constant symbol in the algebraic signature).  Formally, there
exists a function $\mathrm{witness} : C(\alpha) \to \alpha$.
\end{definition}

\begin{theorem}[General paradigm boundary]%
\label{thm:machine_boundary}
Let $C$ be a constant-dependent type class.  If $\alpha$ is empty,
then no instance of $C$ on $\alpha$ can exist.
\end{theorem}

\begin{proof}
Let $i$ be a hypothetical instance of $C$ on $\alpha$.
Since $C$ is constant-dependent, $\mathrm{witness}(i)$ is an
element of $\alpha$.  But $\alpha$ is empty.  Contradiction.
\end{proof}

\noindent The algebraic hierarchy---\texttt{Zero}, \texttt{One},
\texttt{CommSemiring}, \texttt{Semiring}, \texttt{Ring},
\texttt{Field}, and even \texttt{Inhabited}---is registered as
constant-dependent by providing the witness function (typically
extracting $0$ via \texttt{inst.toZero}).  Each registration is a
one-line instance declaration; the impossibility theorem follows
automatically.

To check whether a tactic can cross the paradigm boundary, check
whether its required type class has a \texttt{ConstantDependent}
instance.  If yes, the tactic is locked to inhabited types.  If
no, it is free to cross.

\begin{lstlisting}
-- The boundary classification class
class ConstantDependent
    (C : Type u (*$\to$*) Type v) where
  witness : {(*$\alpha$*) : Type u} (*$\to$*) C (*$\alpha$*) (*$\to$*) (*$\alpha$*)

-- Register the algebraic hierarchy:
instance : ConstantDependent Zero where
  witness inst := @Zero.zero _ inst
instance : ConstantDependent CommSemiring where
  witness inst := @Zero.zero _ inst.toZero
instance : ConstantDependent Field where
  witness inst := @Zero.zero _ inst.toZero
instance : ConstantDependent Inhabited where
  witness inst := @Inhabited.default _ inst

-- ONE theorem covers them all:
theorem constant_dependent_blocks_empty
    {C : Type u (*$\to$*) Type v}
    [ConstantDependent C]
    {(*$\alpha$*) : Type u}
    (hempty : IsEmpty (*$\alpha$*))
    (inst : C (*$\alpha$*)) : False :=
  hempty.false (ConstantDependent.witness inst)

-- Every impossibility is now a one-liner:
example ((*$\alpha$*) : Type*) (h : IsEmpty (*$\alpha$*))
    (i : Ring (*$\alpha$*)) : False :=
  constant_dependent_blocks_empty h i
\end{lstlisting}

\noindent All formalizations---three layers plus the
\texttt{ConstantDependent} hierarchy---compile with Lean~4.29.0
and Mathlib~4.29.0, with zero \texttt{sorry} declarations.
The paradigm boundary discovered empirically in Parts~1--10 is
now a formal, machine-checkable theorem.

\subsection{Connections to foundational ideas}

The paradigm boundary is not an isolated observation about Lean
tactics.  It connects to several foundational ideas in logic,
type theory, and category theory.

\paragraph{Curry--Howard and the boundary as a type-theoretic
phenomenon.}
Under the Curry--Howard
correspondence~\cite{howard1980, curry1958},
proofs are programs and propositions are types.  The paradigm
boundary is, at bottom, a statement about which types can serve as
domains for which programs.  The \texttt{ring} tactic constructs a
\emph{program} (a proof term) that manipulates $0$, $1$, $+$,
$\times$ via the \texttt{CommSemiring} axioms.  This program has
type $a = b$ where $a, b : R$ and $R$ carries \texttt{CommSemiring}.
The Curry--Howard reading is: the \emph{proposition} ``$a = b$ in a
commutative semiring'' can only be \emph{inhabited} (proved) when the
type $R$ itself is inhabited---because the proof term must manipulate
elements of~$R$.  The structural fragment, by contrast, constructs
proof terms that never mention elements of the carrier: \texttt{funext}
quantifies over elements but does not produce them; \texttt{congr}
decomposes equality but does not evaluate subterms.  These proof
terms inhabit their proposition types regardless of whether the
carrier is inhabited.  The paradigm boundary is thus a
\emph{habitation boundary} in the Curry--Howard sense: it separates
propositions whose proofs require witness terms from propositions
whose proofs are purely structural.

\paragraph{Type-class hierarchies across proof assistants.}
The boundary's formal mechanism---\texttt{Zero} providing a
constant symbol that forces nonemptiness---is specific to how
Mathlib organizes its algebraic hierarchy.  In Lean~4's Mathlib,
\texttt{CommSemiring} extends \texttt{Zero} and \texttt{One}
as explicit type-class components, and the \texttt{ring} tactic's
elaborator synthesizes these instances at invocation
time~\cite{mathlib2020, moura2021}.  In
Coq's Mathematical Components library~\cite{mahboubi2013},
the analogous hierarchy uses \emph{packed classes}:
\texttt{GRing.Ring} bundles a zero, a one, and all axioms into a
single structure.  The effect is the same---you cannot form
\texttt{GRing.Ring} over an empty type---but the mechanism differs:
Coq's canonical structures resolve at unification time rather than
via instance search.  In both systems, the paradigm boundary
holds: algebraic decision procedures require an inhabited carrier.
The boundary is not an artifact of Lean's design but a consequence
of the algebraic signature itself: any proof assistant that
formalizes ring theory must, by the axioms, produce a zero element,
and any empty type contradicts this.

\paragraph{Category theory: natural transformations vs.\ internal
algebra.}
The two proof fragments correspond to two levels of categorical
abstraction.  The \emph{structural fragment} operates at the level
of \emph{natural transformations}: functoriality, extensionality,
and composition.  These are \emph{element-free} notions---a natural
transformation $\eta : F \Rightarrow G$ is defined by its components
$\eta_X : F(X) \to G(X)$ for each object~$X$, but the definition
does not require any particular object to exist.  Natural
transformations between functors on the empty category are trivially
well-defined (there are no components to specify).  This is why
the structural fragment works on empty types: it lives in the
\emph{external} categorical world where objects may or may not have
elements.

The \emph{algebraic fragment} operates at the level of
\emph{internal algebraic structures}: commutative semirings, fields,
ordered rings.  These are defined by operations on \emph{elements}
of a carrier---elements that must exist for the operations to be
meaningful.  In categorical terms, an internal ring in a category
$\mathcal{C}$ is an object $R$ equipped with morphisms
$0 : \mathbf{1} \to R$, $+ : R \times R \to R$, etc., where
$\mathbf{1}$ is the terminal object.  The morphism
$0 : \mathbf{1} \to R$ is a \emph{global element}---it picks out a
specific point of~$R$.  If $R$ is empty (initial rather than
terminal), no such morphism exists.

The paradigm boundary is therefore the boundary between
\emph{external categorical reasoning} (natural transformations,
functoriality, universality) and \emph{internal algebraic reasoning}
(elements, operations, equations).  This distinction is well-known
in category theory~\cite{maclane1971, borceux1994}; what the Yanasse
campaign adds is the empirical observation that it manifests as a
\emph{tactic transfer barrier} in formalized mathematics, and the
formal proof that this barrier is logically necessary.

\paragraph{On the apparent triviality of the result.}
A skeptic might object: ``We already know rings are not empty.''
Indeed---but knowing it is intuition.  Formalizing \emph{the exact
moment the logic breaks}---the point where a constant symbol in an
algebraic signature forces the carrier to be nonempty, and the
precise type-class chain through which this constraint propagates
to the tactic elaborator---allows us to build better automated
search tools.  The intuition that ``\texttt{ring} won't work here''
is correct but unhelpful: it does not tell the prover \emph{when}
to stop trying, \emph{what} to extract from the failed attempt, or
\emph{where} to look instead.  The formal boundary does all three.
We have, in effect, measured the friction coefficient of mathematics
itself: the quantitative resistance that the type-class landscape
offers to tactic transfer across area boundaries.

\subsection{Structural density: a metric for proof reusability}

The paradigm boundary suggests a new metric for the formal
verification community.

\begin{definition}[Structural density]%
\label{def:structural_density}
The \emph{structural density} of a proof is the fraction of its
tactic steps that belong to the constant-free (structural)
fragment---steps that use no algebraic constants and require no
type-class instances providing elements.  Formally, if a proof
consists of $n$~tactic steps, of which $s$~are structural
(extensionality, congruence, case splits, function application)
and $a = n - s$ are algebraic (ring normalization, linear
arithmetic, field simplification, numerical evaluation), then
\[
  \rho \;=\; \frac{s}{n}
\]
is the structural density, with $\rho \in [0, 1]$.
\end{definition}

\noindent The Yanasse campaign provides empirical calibration:
structural tactics transfer at ${\sim}$60\% across area boundaries;
algebraic tactics at ${\sim}$10\%.  A proof's structural density
therefore predicts its \emph{transferability}---the probability
that the proof technique can be reused in a distant area of
mathematics.

\paragraph{High structural density ($\rho \to 1$).}
A proof composed almost entirely of extensionality reductions,
case splits, and congruence steps is \emph{reusable}: its technique
can transfer to any target area because it makes no assumptions
about the algebraic structure of the types involved.  Such proofs
are valuable in a library precisely because they are
portable---they work wherever the logical structure matches,
regardless of the algebraic context.

\paragraph{Low structural density ($\rho \to 0$).}
A proof composed almost entirely of \texttt{ring}, \texttt{linarith},
and \texttt{norm\_num} is \emph{brittle}: it works only in the
specific algebraic neighborhood where the required type-class
instances exist.  Such proofs are efficient within their home area
but contribute nothing to cross-area technique transfer.  They are
the mathematical equivalent of highly optimized but non-portable
code.

\paragraph{Implications for library design.}
If a theorem admits two proofs---one with $\rho = 0.8$ and one with
$\rho = 0.3$---the higher-density proof is strictly more valuable
to the library as a \emph{source of transferable techniques}, even
if the lower-density proof is shorter or more ``elegant'' by
traditional metrics.  Structural density could therefore serve as
a secondary quality criterion for proof libraries: among correct
proofs of the same theorem, prefer the one with higher $\rho$,
because it maximizes the theorem's utility as an analogy source.

\paragraph{Computed values across Mathlib.}
We classified each of Mathlib's 209{,}519~tactic invocations as
structural ($\rho$-contribution~$= 1$), algebraic ($= 0$), or
mixed ($= 0.1$--$0.9$), using the empirical transfer rates from
the campaign as the classification guide.
Table~\ref{tab:structural_density} shows $\rho$ for all
27~areas.  The Mathlib-wide average is
$\bar{\rho} = 0.790$: about 79\% of all tactic steps in
Mathlib are structural.

\begin{table}[ht]
\centering\small
\caption{Structural density $\rho$ by Mathlib area, computed over
  209{,}519 tactic invocations.  Higher $\rho$ = more structural
  (reusable); lower $\rho$ = more algebraic (brittle).}
\label{tab:structural_density}
\begin{tabular}{@{}lrc@{\quad}lrc@{}}
\toprule
Area & Tactics & $\rho$ & Area & Tactics & $\rho$ \\
\midrule
FieldTheory     &  2{,}773 & .840 &
  Topology       & 13{,}304 & .792 \\
Control         &    246 & .838 &
  AlgGeometry    &  5{,}456 & .789 \\
SetTheory       &  2{,}071 & .835 &
  Combinatorics  &  7{,}474 & .789 \\
Dynamics        &    493 & .824 &
  AlgTopology    &  2{,}288 & .782 \\
GroupTheory     &  5{,}643 & .819 &
  MeasureTheory  & 17{,}489 & .781 \\
Logic           &    656 & .817 &
  Computability  &  2{,}344 & .769 \\
ModelTheory     &  1{,}079 & .816 &
  Tactic         &    362 & .762 \\
Condensed       &    187 & .816 &
  Analysis       & 30{,}328 & .756 \\
Algebra         & 27{,}721 & .814 &
  InfoTheory     &    240 & .752 \\
Order           &  4{,}817 & .810 &
  Probability    &  6{,}183 & .752 \\
RepTheory       &    669 & .807 &
  NumberTheory   &  8{,}577 & .742 \\
RingTheory      & 20{,}334 & .805 &
  & & \\
LinearAlgebra   & 11{,}300 & .805 &
  \textbf{Mathlib-wide} & \textbf{209{,}519} & \textbf{.790} \\
Data            & 12{,}652 & .796 &
  & & \\
Geometry        &  6{,}732 & .795 &
  CatTheory      & 18{,}101 & .794 \\
\bottomrule
\end{tabular}
\end{table}

\paragraph{What $\rho$ does and does not predict.}
Structural density does \emph{not} directly predict transfer
success rate.  Probability ($\rho = 0.752$) and Number Theory
($\rho = 0.742$) have nearly identical structural densities but
4.6$\times$ different transfer rates (46\% vs.\ 10\%).  This is
because $\rho$ measures a proof's \emph{style}---the mix of
structural and algebraic steps---not its \emph{transferability},
which depends on the two-factor model of
Section~\ref{sec:model} (source tactic profile $\times$ target
type-class landscape).

What $\rho$ \emph{does} predict is \emph{reusability}: a high-$\rho$
proof contains more structural skeleton that can survive transfer.
Field Theory ($\rho = 0.840$) and Set Theory ($\rho = 0.835$)
have the most portable proofs; Number Theory ($\rho = 0.742$)
and Probability ($\rho = 0.752$) have the least.  The
$\rho$~values, together with the transfer potential matrices of
Section~\ref{sec:map}, form a complete \emph{reusability map} of
Mathlib: the matrices say where to look for analogues; $\rho$
says how much of what you find will survive the transfer.

%======================================================================
\section{Conclusion}\label{sec:conclusion}
%======================================================================

The initial Yanasse campaign---100~schema transfer attempts across
10~area pairs, producing 35~Lean-verified new proofs---demonstrates
that a domain-independent relational analogy engine can systematically
discover cross-domain connections between areas of formalized
mathematics.  The engine sees the ``shape'' of proofs without
understanding their content; when that shape aligns with type-class
compatibility, new proofs emerge.

The campaign's most important contribution is not the 35~new proofs
themselves---each is a new proof of an existing theorem, not a new
theorem in the mathematical sense---but the \emph{structural map} of
transferability across Mathlib.  This map reveals:

\begin{enumerate}
\item A \textbf{two-factor model}: transfer success depends on
  source tactic profile $\times$ target type-class landscape.

\item An \textbf{element/morphism paradigm boundary}: algebraic
  decision procedures cannot cross from element-level to
  morphism-level proof domains, except where element-level
  sub-layers exist.

\item A \textbf{tactic taxonomy}: structural combinators transfer
  broadly (5+/7~targets); algebraic procedures transfer narrowly
  (1/3~targets, and only with type-class compatibility);
  meta-tactics transfer independently of QAP scores.

\item A \textbf{complementarity principle}: different source areas
  access non-overlapping layers of the same target, so multi-source
  campaigns achieve broader coverage.
\end{enumerate}

These findings directly inform the Deep Vision theorem prover:
prioritize Probability as a general-purpose analogy source; add
type-class pre-filtering to the QAP matcher; weight analogy
suggestions by the two-factor model; and maintain separate libraries
for structural and algebraic tactic families.

The Yanasse method is not limited to the two source areas tested.
The pipeline is fully automated (GPU-accelerated matching +
AI-driven semantic adaptation + Lean compilation), and the
structural map suggests that new source and target areas would each
contribute unique transfer vectors.

\subsection*{The next ten pairs}

The initial campaign processed the top~10 pairs from
\texttt{area\_pair\_potential.parquet}, ranked by the naive model
(sum of top-10 schema transfer gaps weighted by source-count).
Section~\ref{sec:map} introduced a transferability-weighted model
that incorporates the campaign's empirical transfer rates per tactic
class.  The next 10~pairs are selected by the weighted model
(Table~\ref{tab:weighted_matrix}), which reshuffles the rankings
dramatically---Number Theory pairs drop out entirely, and three new
source areas appear:

\medskip
{\small
\begin{center}
\begin{tabular}{@{}rlccp{4.8cm}@{}}
\toprule
Rank & Pair & Pot.\ & Cand.\ & Prediction \\
\midrule
1 & Prob $\to$ Computability & 68 & 19
  & High: structural combinators on decidability proofs \\
2 & Prob $\to$ AlgTopology & 68 & 20
  & High: structural combinators + shared topology \\
3 & Prob $\to$ GroupTheory & 57 & 16
  & High: structural combinators on algebraic targets \\
4 & AlgGeom $\to$ Computability & 56 & 30
  & Unknown: first test of AlgGeom as source \\
5 & AlgGeom $\to$ Dynamics & 55 & 30
  & Unknown: scheme-level tactics on mixed target \\
6 & Prob $\to$ CategoryTheory & 55 & 17
  & Moderate: structural tactics on categorical proofs \\
7 & RepTheory $\to$ InfoTheory & 54 & 18
  & Unknown: first test of RepTheory as source \\
8 & Prob $\to$ LinearAlgebra & 54 & 15
  & High: structural combinators + algebraic structure \\
9 & AlgGeom $\to$ InfoTheory & 53 & 33
  & Unknown: categorical + algebraic hybrid \\
10 & AlgGeom $\to$ Condensed & 53 & 33
  & Unknown: scheme-level tactics on categorical target \\
\bottomrule
\end{tabular}
\end{center}
}

\noindent The weighted model makes several testable predictions.
Probability~$\to$~Computability (rank~1) and
Probability~$\to$~AlgebraicTopology (rank~2) should match or exceed
Probability's 46\% campaign average, since their candidates are
dominated by structural combinators (\texttt{ext1},
\texttt{by\_cases}, \texttt{any\_goals}).
Algebraic Geometry (ranks~4, 5, 9, 10) is the most interesting
unknown: its tactic profile mixes categorical DSL-style tactics
with scheme-level algebraic computation, and the campaign offers
no prior data on how this hybrid transfers.
Representation Theory~$\to$~Information Theory (rank~7) tests
whether a \emph{target} area from the initial campaign can serve
as a \emph{source}---a role reversal that the two-factor model
predicts should work if RepTheory's distinctive tactics (chain
complex decomposition, module extensionality) have structural
analogs in InfoTheory's proof architecture.

These 10~pairs are the next phase of the Yanasse project.  Combined
with the initial campaign, they will cover 4~source areas
(Probability, Algebraic Geometry, Representation Theory, plus
continued Number Theory analysis),
9~target areas, and 200~schema transfer attempts---doubling the
empirical base for the structural map of Mathlib transferability.

\subsection*{A strange loop}

The Yanasse project, viewed whole, forms what
Hofstadter~\cite{hofstadter_go_1999, hofstadter2007loop,
hofstadter_surfaces_2013}
calls a \emph{strange loop}: a system that moves through levels of
abstraction and arrives back at itself.

Deep Vision was built to understand chess---to see that a pawn
defended by its king and attacked by a knight plays ``the same
role'' as a rook defended by its queen and attacked by a bishop.
The same relational matching engine, with only the relation
extractor swapped, was applied to Lean~4 proof states: hypotheses,
goals, and types replaced pieces; rewrite, apply, and head-match
replaced attack, defense, and blocking.  This engine, running at
scale across 217{,}133 proof states (the Yanasse campaign),
discovered empirically that certain tactic transfers are
impossible---20~attempts against Condensed Mathematics and
Representation Theory, 0~successes, every failure tracing to a
missing type-class instance.  We called this the paradigm boundary.

We then formalized the paradigm boundary as a Lean theorem and fed
it back to the proving loop---the same relational analogy engine.
The engine matched the theorem's proof state against its library,
found structurally similar proofs in measure theory, combinatorics,
and analysis (completely unrelated mathematics), and extracted the
right strategy: assume the algebraic instance, extract a zero
element, contradict emptiness.  The analogy is absurd at the
content level---what does a graph-coloring argument have to do
with the non-existence of commutative semirings on empty
types?---but structurally valid: both involve deriving a
contradiction from an existence assumption.  The system perceived
the shared essence across superficially different situations, which
is Hofstadter's definition of analogy.

But the engine could not compile the final proof.  The obstacle was
type-class threading: after introducing the hypothetical
\texttt{CommSemiring} instance, Lean's synthesizer did not
recognize it as available for resolving \texttt{(0\,:\,$\alpha$)}.
The proof required an explicit annotation---\texttt{@Zero.zero}\,$\alpha$\,\texttt{inst.toZero}---to
route through the locally introduced hypothesis.

This failure is the strange loop.  The system's limitation in
\emph{proving} the paradigm boundary theorem is an instance of the
paradigm boundary \emph{itself}: the proving loop's adaptation
layer could not thread a type-class instance across a local
context boundary, which is precisely the kind of type-class
incompatibility that the theorem formalizes.  The proof of the
boundary ran into the boundary.

The full loop: chess $\to$ proof states $\to$ tactic transfer $\to$
paradigm boundary $\to$ formal theorem $\to$ feed back to the
matcher $\to$ the matcher's failure illuminates why the boundary
exists.  Each level of abstraction leads to the next, and the last
level reflects back on the first.  The domain-independent matching
engine that sees ``this pawn plays the same role as that rook''
also sees ``this impossibility theorem has the same shape as that
measure-zero argument''---and its inability to close the proof is
itself a data point about the structure of mathematical knowledge.

\section*{Acknowledgments}

Named after Horacio Hideki Yanasse.  The Deep Vision matching engine
was originally developed for chess cognition research; its application
to formal mathematics was enabled by the domain-independent design of
\texttt{deep\_vision\_lib.py}.  All computations were performed on a
MacBook Air (Apple M-series, 32\,GB unified memory) using PyTorch~2.8
on the MPS backend.

%======================================================================
\begin{thebibliography}{99}
%======================================================================

\bibitem{howard1980}
W.~A.~Howard,
``The formulae-as-types notion of construction,''
in \emph{To H.~B.~Curry: Essays on Combinatory Logic, Lambda
Calculus and Formalism} (J.~P.~Seldin and J.~R.~Hindley, eds.),
pp.~479--490, Academic Press, 1980.

\bibitem{curry1958}
H.~B.~Curry, R.~Feys, and W.~Craig,
\emph{Combinatory Logic}, vol.~1.
Amsterdam: North-Holland, 1958.

\bibitem{mathlib2020}
The Mathlib Community,
``The Lean mathematical library,''
in \emph{Proceedings of the 9th ACM SIGPLAN International Conference
on Certified Programs and Proofs (CPP 2020)}, pp.~367--381, 2020.

\bibitem{moura2021}
L.~de~Moura, S.~Kong, J.~Avigad, F.~van~Doorn, and M.~von~Raumer,
``The Lean theorem prover (system description),''
in \emph{Automated Deduction --- CADE-25}
(A.~P.~Felty and A.~Middeldorp, eds.),
Lecture Notes in Computer Science, vol.~9195, pp.~378--388,
Springer, 2015.

\bibitem{mahboubi2013}
A.~Mahboubi and E.~Tassi,
``Canonical structures for the working Coq user,''
in \emph{Interactive Theorem Proving (ITP 2013)}
(S.~Blazy, C.~Paulin-Mohring, and D.~Pichardie, eds.),
Lecture Notes in Computer Science, vol.~7998, pp.~19--34,
Springer, 2013.

\bibitem{maclane1971}
S.~Mac~Lane,
\emph{Categories for the Working Mathematician}.
New York: Springer-Verlag, 1971.

\bibitem{borceux1994}
F.~Borceux,
\emph{Handbook of Categorical Algebra}, 3~vols.
Cambridge: Cambridge University Press, 1994.

\bibitem{yanasse_part1}
A.~Linhares,
``Yanasse: Finding new proofs from Deep Vision's analogies --- Part~1:
Probability Theory $\to$ Representation Theory,''
ARGO LABS Report for Distribution, April 2026.

\bibitem{yanasse_part2}
A.~Linhares,
``Yanasse: Finding new proofs from Deep Vision's analogies --- Part~2:
Probability Theory $\to$ Set Theory,''
ARGO LABS Report for Distribution, April 2026.

\bibitem{yanasse_part3}
A.~Linhares,
``Yanasse: Finding new proofs from Deep Vision's analogies --- Part~3:
Probability Theory $\to$ Condensed Mathematics,''
ARGO LABS Report for Distribution, April 2026.

\bibitem{yanasse_part4}
A.~Linhares,
``Yanasse: Finding new proofs from Deep Vision's analogies --- Part~4:
Probability Theory $\to$ Field Theory,''
ARGO LABS Report for Distribution, April 2026.

\bibitem{yanasse_part5}
A.~Linhares,
``Yanasse: Finding new proofs from Deep Vision's analogies --- Part~5:
Probability Theory $\to$ Dynamics,''
ARGO LABS Report for Distribution, April 2026.

\bibitem{yanasse_part6}
A.~Linhares,
``Yanasse: Finding new proofs from Deep Vision's analogies --- Part~6:
Probability Theory $\to$ Model Theory,''
ARGO LABS Report for Distribution, April 2026.

\bibitem{yanasse_part7}
A.~Linhares,
``Yanasse: Finding new proofs from Deep Vision's analogies --- Part~7:
Probability Theory $\to$ Information Theory,''
ARGO LABS Report for Distribution, April 2026.

\bibitem{yanasse_part8}
A.~Linhares,
``Yanasse: Finding new proofs from Deep Vision's analogies --- Part~8:
Number Theory $\to$ Condensed Mathematics,''
ARGO LABS Report for Distribution, April 2026.

\bibitem{yanasse_part9}
A.~Linhares,
``Yanasse: Finding new proofs from Deep Vision's analogies --- Part~9:
Number Theory $\to$ Representation Theory,''
ARGO LABS Report for Distribution, April 2026.

\bibitem{yanasse_part10}
A.~Linhares,
``Yanasse: Finding new proofs from Deep Vision's analogies --- Part~10:
Number Theory $\to$ Dynamics,''
ARGO LABS Report for Distribution, April 2026.

\bibitem{linhares2024deepvision}
A.~Linhares,
``Deep Vision,''
ARGO LABS Internal Report, 2026.

\bibitem{hofstadter_fluid_1995}
D.~R.~Hofstadter and the Fluid Analogies Research Group,
\emph{Fluid Concepts \& Creative Analogies: Computer Models of the
Fundamental Mechanisms of Thought}. New York: Basic Books, 1995.

\bibitem{hofstadter_go_1999}
D.~R.~Hofstadter,
\emph{G\"odel, Escher, Bach: An Eternal Golden Braid},
20th anniversary ed.\ New York: Basic Books, 1999.

\bibitem{hofstadter_surfaces_2013}
D.~R.~Hofstadter and E.~Sander,
\emph{Surfaces and Essences: Analogy as the Fuel and Fire of
Thinking}. New York: Basic Books, 2013.

\bibitem{linhares_active_2005}
A.~Linhares,
``An active symbols theory of chess intuition,''
\emph{Minds and Machines}, vol.~15, pp.~131--151, 2005.

\bibitem{linhares_understanding_2007}
A.~Linhares and P.~Brum,
``Understanding our understanding of strategic scenarios: What role
do chunks play?''
\emph{Cognitive Science}, vol.~31, no.~6, pp.~989--1007, 2007.

\bibitem{linhares_questioning_2010}
A.~Linhares and A.~E.~T.~A.~Freitas,
``Questioning Chase and Simon's (1973) `Perception in Chess',''
\emph{New Ideas in Psychology}, vol.~28, no.~1, pp.~64--78, 2010.

\bibitem{linhares_entanglement_2012}
A.~Linhares, A.~E.~T.~A.~Freitas, A.~Mendes, and J.~S.~Silva,
``Entanglement of perception and reasoning in the combinatorial
game of chess,''
\emph{Cognitive Systems Research}, vol.~13, no.~1, pp.~72--86, 2012.

\bibitem{linhares_what_2013}
A.~Linhares and D.~M.~Chada,
``What is the nature of the mind's pattern-recognition process?''
\emph{New Ideas in Psychology}, vol.~31, no.~2, pp.~108--121, 2013.

\bibitem{linhares_emergence_2014}
A.~Linhares,
``The emergence of choice: Decision-making and strategic thinking
through analogies,''
\emph{Information Sciences}, vol.~259, pp.~36--56, 2014.

\bibitem{linhares_deep_vision}
A.~Linhares,
``Deep Vision: Seeing the essence of proofs through analogy,''
ARGO LABS Internal Report, 2026.

\bibitem{linhares_wolstenholme_2026}
A.~Linhares,
``Deep Vision: A formal proof of Wolstenholme's theorem in Lean~4,''
ARGO LABS Report for Distribution, 2026.

\bibitem{linhares_sketch_2026}
A.~Linhares,
``A sketch of Deep Vision's study of mathematics,''
ARGO LABS Report for Distribution, 2026.

% === New references (post-2010, from target journals/arXiv) ===

\bibitem{hofstadter2007loop}
D.~R.~Hofstadter,
\emph{I Am a Strange Loop}.
New York: Basic Books, 2007.

\bibitem{polu2020gptf}
S.~Polu and I.~Sutskever,
``Generative language modeling for automated theorem proving,''
arXiv preprint 2009.03393, 2020.

\bibitem{han2022pact}
J.~M.~Han, J.~Rute, Y.~Wu, E.~W.~Ayers, and S.~Polu,
``Proof artifact co-training for theorem proving with language
models,''
in \emph{Proceedings of the International Conference on Learning
Representations (ICLR)}, 2022.

\bibitem{bansal2019holist}
K.~Bansal, S.~Loos, M.~Rabe, C.~Szegedy, and S.~Wilcox,
``HOList: An environment for machine learning of higher-order
theorem proving,''
in \emph{Proceedings of the 36th International Conference on Machine
Learning (ICML)}, vol.~97, pp.~454--463, PMLR, 2019.

\bibitem{yang2024leandojo}
K.~Yang, A.~M.~Swope, A.~Gu, R.~Chalamala, P.~Song, S.~Yu,
S.~Godil, R.~Prenger, and A.~Anandkumar,
``LeanDojo: Theorem proving with retrieval-augmented language
models,''
in \emph{Advances in Neural Information Processing Systems~36
(NeurIPS)}, 2023.

\bibitem{song2024leancopilot}
P.~Song, K.~Yang, and A.~Anandkumar,
``Lean Copilot: Large language models as copilots for theorem
proving in Lean,''
in \emph{Proceedings of Machine Learning Research}, vol.~288,
pp.~1--26, 2025.

\bibitem{xin2024deepseekprover}
H.~Xin, D.~Guo, Z.~Shao, Z.~Ren, Q.~Zhu, B.~Liu, C.~Ruan,
W.~Li, and X.~Liang,
``DeepSeek-Prover: Advancing theorem proving in LLMs through
large-scale synthetic data,''
arXiv preprint 2405.14333, 2024.

\bibitem{wang2024legoprover}
H.~Wang, H.~Xin, C.~Zheng, L.~Li, Z.~Liu, Q.~Cao, Y.~Huang,
J.~Xiong, H.~Shi, E.~Xie, J.~Yin, Z.~Li, and H.~Liang,
``LEGO-Prover: Neural theorem proving with growing libraries,''
in \emph{Proceedings of the International Conference on Learning
Representations (ICLR)}, 2024.

\bibitem{trinh2024alphageometry}
T.~H.~Trinh, Y.~Wu, Q.~V.~Le, H.~He, and T.~Luong,
``Solving olympiad geometry without human demonstrations,''
\emph{Nature}, vol.~625, pp.~476--482, 2024.

\bibitem{alphaproof2025nature}
AlphaProof Team (Google DeepMind),
``Olympiad-level formal mathematical reasoning with reinforcement
learning,''
\emph{Nature}, 2025.

\bibitem{li2024survey}
Z.~Li, J.~Sun, L.~Murphy, Q.~Su, Z.~Li, X.~Zhang, K.~Yang,
and X.~Si,
``A survey on deep learning for theorem proving,''
in \emph{Proceedings of the Conference on Language Modeling
(COLM)}, 2024.

\bibitem{kaliszyk2014flyspeck}
C.~Kaliszyk and J.~Urban,
``Learning-assisted automated reasoning with Flyspeck,''
\emph{Journal of Automated Reasoning}, vol.~53, pp.~173--213,
2014.

\bibitem{blanchette2016hammering}
J.~C.~Blanchette, C.~Kaliszyk, L.~C.~Paulson, and J.~Urban,
``Hammering towards QED,''
\emph{Journal of Formalized Reasoning}, vol.~9, no.~1,
pp.~101--148, 2016.

\bibitem{gauthier2021tactictoe}
T.~Gauthier, C.~Kaliszyk, J.~Urban, R.~Kumar, and M.~Norrish,
``TacticToe: Learning to prove with tactics,''
\emph{Journal of Automated Reasoning}, vol.~65, pp.~257--286,
2021.

\bibitem{mikula2024magnushammer}
M.~Mikula, S.~Tworkowski, S.~Antoniak, B.~Piotrowski,
A.~Q.~Jiang, J.~P.~Zhou, C.~Szegedy, {\L}.~Kuci\'nski,
P.~Mi{\l}o{\'s}, and Y.~Wu,
``Magnushammer: A transformer-based approach to premise selection,''
in \emph{Proceedings of the International Conference on Learning
Representations (ICLR)}, 2024.

\bibitem{gauthier2019aligning}
T.~Gauthier and C.~Kaliszyk,
``Aligning concepts across proof assistant libraries,''
\emph{Journal of Symbolic Computation}, vol.~90, pp.~89--123,
2019.

\bibitem{ringer2021proofrepair}
T.~Ringer, R.~Porter, N.~Yazdani, J.~Leo, and D.~Grossman,
``Proof repair across type equivalences,''
in \emph{Proceedings of the 42nd ACM SIGPLAN International
Conference on Programming Language Design and Implementation
(PLDI)}, pp.~112--128, 2021.

\bibitem{gentner1983structure}
D.~Gentner,
``Structure-mapping: A theoretical framework for analogy,''
\emph{Cognitive Science}, vol.~7, no.~2, pp.~155--170, 1983.

\bibitem{mitchell2021analogy}
M.~Mitchell,
``Abstraction and analogy-making in artificial intelligence,''
\emph{Annals of the New York Academy of Sciences}, vol.~1505,
pp.~79--101, 2021.

\bibitem{moura2021lean4}
L.~de~Moura and S.~Ullrich,
``The Lean~4 theorem prover and programming language,''
in \emph{Automated Deduction---CADE~28}
(A.~Platzer and G.~Sutcliffe, eds.),
Lecture Notes in Computer Science, vol.~12699, pp.~625--635,
Springer, 2021.

\bibitem{buzzard2020perfectoid}
K.~Buzzard, J.~Commelin, and P.~Massot,
``Formalising perfectoid spaces,''
in \emph{Proceedings of the 9th ACM SIGPLAN International
Conference on Certified Programs and Proofs (CPP)},
pp.~299--312, 2020.

\bibitem{scholze2022lte}
P.~Scholze and the Lean Community,
``Completion of the Liquid Tensor Experiment,''
Lean Community Blog, 2022.
% \url{https://leanprover-community.github.io/blog/posts/lte-final/}

\bibitem{avigad2024formalturn}
J.~Avigad,
``Mathematics and the formal turn,''
\emph{Bulletin of the American Mathematical Society}, vol.~61,
no.~2, pp.~225--240, 2024.

\bibitem{zheng2022minif2f}
K.~Zheng, J.~M.~Han, and S.~Polu,
``MiniF2F: A cross-system benchmark for formal Olympiad-level
mathematics,''
in \emph{Proceedings of the International Conference on Learning
Representations (ICLR)}, 2022.

\bibitem{wu2022autoformalization}
Y.~Wu, A.~Q.~Jiang, W.~Li, M.~Rabe, C.~Staats, M.~Jamnik,
and C.~Szegedy,
``Autoformalization with large language models,''
in \emph{Advances in Neural Information Processing Systems~35
(NeurIPS)}, 2022.

\bibitem{jiang2023dsp}
A.~Q.~Jiang, S.~Welleck, J.~P.~Zhou, W.~Li, J.~Liu,
M.~Jamnik, T.~Lacroix, Y.~Wu, and G.~Lample,
``Draft, sketch, and prove: Guiding formal theorem provers with
informal proofs,''
in \emph{Proceedings of the International Conference on Learning
Representations (ICLR)}, 2023.

\bibitem{szegedy2020autoformalization}
C.~Szegedy,
``A promising path towards autoformalization and general artificial
intelligence,''
in \emph{Intelligent Computer Mathematics (CICM 2020)}
(C.~Benzm\"uller and B.~Miller, eds.),
Lecture Notes in Computer Science, vol.~12236, pp.~3--20,
Springer, 2020.

\bibitem{babai2016quasipoly}
L.~Babai,
``Graph isomorphism in quasipolynomial time,''
in \emph{Proceedings of the 48th Annual ACM Symposium on Theory
of Computing (STOC)}, pp.~684--697, 2016.

\end{thebibliography}

\end{document}
